%% rset.applications.tex
%% Applications: the Pitts--Gabbay nominal theory, founded on the reflective
%% construction via the FM model of Section~\ref{sec:fmmodel}.
%% Drop in via %% rset.applications.tex
%% Applications: the Pitts--Gabbay nominal theory, founded on the reflective
%% construction via the FM model of Section~\ref{sec:fmmodel}.
%% Drop in via %% rset.applications.tex
%% Applications: the Pitts--Gabbay nominal theory, founded on the reflective
%% construction via the FM model of Section~\ref{sec:fmmodel}.
%% Drop in via %% rset.applications.tex
%% Applications: the Pitts--Gabbay nominal theory, founded on the reflective
%% construction via the FM model of Section~\ref{sec:fmmodel}.
%% Drop in via \input{rset.applications} (add the hook to rset.tex after
%% \input{rset.mainthm}).
\input{rset.colors}   % red/black colour layer + FM symbols (providecommand-guarded)

\section{Applications: nominal techniques founded on reflection}
\label{sec:applications}

The original motivation for $\mathsf{FM}$ in computer science was Pitts and
Gabbay's \emph{nominal} treatment of names and binding
\cite{DBLP:journals/fac/GabbayP02}. That treatment is founded squarely on
$\mathsf{FM}$: names are the atoms, $\alpha$-equivalence is an orbit relation
under the permutation action, and the freshness machinery is exactly support and
the $\new$-quantifier. Since Section~\ref{sec:fmmodel} realises $\mathsf{FM}$
inside the reflective universe, the entire nominal edifice descends with it. The
names that nominal techniques take as primitive become, here, black sets built
from the empty set; the reflective construction \emph{founds} the nominal theory
rather than presupposing its atoms.

\subsection{Nominal sets, internally}
\label{sec:applications:nom}

A \emph{nominal set} is a set carrying a $\Sym(\Atoms)$-action in which every
element is finitely supported; nominal sets and equivariant maps form the
category $\mathbf{Nom}$, equivalent to the Schanuel topos. The model
$\mathcal{M}$ of Theorem~\ref{thm:rep} is, on the nose, a nominal set: the action
is that of Section~\ref{sec:fmmodel:action}, and finite support is
Corollary~\ref{cor:hfs}. Every construction nominal practitioners rely on ---
the equivariance principle (anything definable from equivariant data is
equivariant), least supports, and the some/any law for $\new$ --- is therefore
available inside the reflective universe, with the colour discipline of
Section~\ref{sec:fmmodel} keeping the bookkeeping honest: names live in the
{\bk black} copy, the sets that bind them in the {\rd red} copy.

\subsection{Abstract syntax with binding}
\label{sec:applications:syntax}

The headline application of \cite{DBLP:journals/fac/GabbayP02} is a uniform
account of inductively-defined syntax \emph{modulo} $\alpha$-equivalence. Its key
device is \emph{name abstraction}: for a nominal set $X$,
\[
  [\Atoms]X \;=\; (\Atoms\times X)\big/\!\sim,
  \qquad
  \langle a\rangle x \sim \langle b\rangle y
  \;\;\Longleftrightarrow\;\;
  \new c.\ (a\,c)\cdot x = (b\,c)\cdot y,
\]
where $\langle a\rangle x$ is the class of $(a,x)$ and $(a\,c)$ is a
transposition. Both the action and $\new$ are present in $\mathcal{M}$, so
$[\Atoms]X$ is definable there, and it has the property that makes it a model of
binding: $\supp(\langle a\rangle x) = \supp(x)\setminus\{a\}$, i.e. abstraction
\emph{discharges} the bound name from the support, precisely as a binder removes
a variable from the free-variable set.

Inductive syntax is then obtained as usual. Object-level $\lambda$-terms, for
instance, are the least solution of
\[
  \Lambda \;\cong\; \Atoms \;+\; \Lambda\times\Lambda \;+\; [\Atoms]\Lambda
  \qquad(\text{variable, application, abstraction}),
\]
with $\alpha$-equivalence built into the third summand. Structural recursion and
induction come with the \emph{freshness condition for binders}, justified by the
some/any law: a definition or proof that treats a bound name as fresh for its
parameters is sound because, $\Atoms$ being infinite, ``some fresh name'' and
``every fresh name'' agree (Corollary~\ref{cor:hfs}). This is exactly the
reasoning principle nominal users invoke when they ``choose a fresh variable''.

The risk accounting carries through and sharpens the point. Each individual
term --- each element of $\Lambda$ --- is a hereditarily finite nominal set, so
it already lives in the $\mathrm{HF}$ model $\mathcal{M}$ of
Theorem~\ref{thm:rep}, founded on the empty set alone: a binding tree of any
fixed shape costs no infinite risk. The completed datatype $\Lambda$, an
infinite equivariant set, is one of the finitely-supported infinities of
Section~\ref{sec:fmmodel:infinitary}, and so costs exactly one $\omega$ --- black
infinity --- and still not a single primitive name. Where Pitts and Gabbay begin
with an infinite set of atoms given by fiat, here the names are manufactured, the
$\alpha$-structure is an orbit of the $\mathrm{swap}$ action, and the whole
apparatus is refinanced down to one empty set and one successor.

\subsection{Nominal algebraic theories}
\label{sec:applications:alg}

Beyond syntax, equational reasoning in the presence of names --- nominal
equational logic and the nominal Lawvere theories of Clouston
\cite{DBLP:journals/jcss/Clouston14} --- takes its semantics in $\mathbf{Nom}$.
As $\mathcal{M}$ furnishes $\mathbf{Nom}$-structure internally, any such theory
(name-restriction, scope extrusion, binding signatures and their equational
laws) can be interpreted in the reflective universe. The reflective setting adds
one structural feature these accounts do not have: by
Corollary~\ref{cor:two} there are \emph{two} interlocking name-universes, the red
and the black, each supplying the other's names. Reading red as player and black
as opponent recovers the game-semantic picture of the \emph{Intuitions}: a name
on one side is a strategy on the other, and a nominal theory and its dual coexist
in a single reflective universe.

\subsection{Summary}
\label{sec:applications:summary}

The chain is short and worth stating plainly. Nominal techniques rest on
$\mathsf{FM}$; $\mathsf{FM}$ embeds in the reflective universe
(Section~\ref{sec:fmmodel}); the reflective universe is built from the empty set.
Composing, the Pitts--Gabbay theory of names and binding is founded on two
recursively intertwined copies of ordinary set theory, with its primitive atoms
replaced by constructed black sets and its infinite risk replaced by the finite
risk of $\varnothing$ --- exactly the reconciliation announced in the
introduction, now delivered for the application that made $\mathsf{FM}$ matter to
computer science in the first place.
 (add the hook to rset.tex after
%% %% rset.mainthm.tex
%% A model of FM set theory inside the reflective theory of Sections 2--3,
%% with the two copies colour-coded (red = player, black = opponent).
%% Drop in via \input{rset.mainthm} (the hook already exists in rset.tex).
\input{rset.colors}   % red/black colour layer + FM symbols (providecommand-guarded)

\section{A model of $\mathsf{FM}$ in the reflective universe}
\label{sec:fmmodel}

We can now collect the dividend of the construction \emph{Tying the second
recursive knot} above. The reflective theory was assembled from the empty set
alone, yet --- we claim --- it already contains a model of $\mathsf{FM}$, and in
fact two of them, interlocked. The atoms such a model needs are not postulated;
they are \emph{supplied by the dual copy}. Throughout we fix the two tied
theories introduced above, the red copy $\rSetX$ and the black copy $\bSetX$,
with
\[
  \rX \;=\; \bSetX, \qquad \bX \;=\; \rSetX,
\]
so that the admissible atoms of the red theory are exactly the black sets, and
\emph{vice versa}. We write the red copy in {\rd red} and the black copy in
{\bk black} throughout, so that a red comprehension $\rbr{\cdots}$ and a black
comprehension $\bbr{\cdots}$, or the red and black membership relations
$\rin$ and $\bin$, are told apart at a glance.

\subsection{Atoms out of the empty set}
\label{sec:fmmodel:atoms}

Inside the black copy, build the von Neumann numerals from the black empty set,
\[
  \bk{0} \;=\; \bemp,
  \qquad
  \bk{n+1} \;=\; \bk{n} \bcup \bbr{\bk{n}},
\]
each a black, hereditarily finite set obtained, ultimately, from $\bemp$. By
black extensionality (the equations of the formal theory above) these are
pairwise distinct: $\bk{n} \neq \bk{m}$ whenever $n\neq m$. Put
\[
  \Atoms \;:=\; \{\, \bk{n} \;:\; n\in\omega \,\}\ \subseteq\ \bSetX .
\]

Two observations make $\Atoms$ behave \emph{as a set of atoms} for the red
theory.
\begin{itemize}
\item \textbf{Admissibility.} By the second recursive knot $\rX=\bSetX$, every
  black set is an admissible red atom; in particular each $\bk{n}$ is one.
\item \textbf{Opacity.} The red membership relation $\rin$ cannot see into a
  black set (the \emph{Intuitions} above). Hence, \emph{as far as the red theory
  can tell}, an element $\bk{n}$ of $\Atoms$ has no internal structure
  whatsoever: the only facts the red theory may assert about $\bk{n}$ and
  $\bk{m}$ are $\bk{n} = \bk{m}$ or $\bk{n} \neq \bk{m}$, and those are settled,
  out of sight, by black equality.
\end{itemize}
Thus $\Atoms$ presents to the red theory precisely as $\mathsf{FM}$'s set of
atoms does: a countable family of mutually distinct, structureless objects that
may sit inside sets but admit no membership query of their own.

\begin{remark}[Where the risk went]
\label{rem:risk}
This is the sense, promised in the introduction, in which atoms cost no infinite
risk. $\mathsf{ZFA}$ and $\mathsf{FM}$ \emph{postulate} an infinite set of
primitive atoms; the postulate is the risk. Here nothing is postulated beyond
$\bemp$. The atoms are black hereditarily finite sets, and the stream $S[X]$ of
the formal theory hands them out lazily, $\mathrm{take}(\Sigma,n)$ at a time, so
that any red set ever built mentions only finitely many of them. The supply is
\emph{potentially} infinite, never an actual completed infinity: the commitment
is the black empty set together with the same successor step that underwrites
$\omega$, and not one primitive atom more.
\end{remark}

\subsection{The permutation action}
\label{sec:fmmodel:action}

The stream operation $\mathrm{swap}(S[X],\rho)$ of the formal theory applies a
finite permutation $\rho=\{\,\bk{j_1}/\bk{i_1}:\dots:\bk{j_k}/\bk{i_k}\,\}$ to
the ordering of atoms; these finite permutations generate the full group
$\Sym(\Atoms)$. A permutation $\pi\in\Sym(\Atoms)$ acts on red sets by
red-membership recursion exactly as in $\mathsf{FM}$,
\[
  \pi\cdot a = \pi(a)\ (a\in\Atoms),
  \qquad
  \pi\cdot S = \rd{\{}\, \pi\cdot e \;:\; e \rin S \,\rd{\}}.
\]
Because $\pi$ moves only atoms and never the black interior of an atom, and
because the red operations $\remp,\ \rcup,\ \rcap,\ \rsm$ and
$\mathrm{embrace}_i$ are visibly equivariant (a renaming commutes with each), the
action is by automorphisms of the red theory and respects the equations of the
formal theory. We have, verbatim, the $\Sym(\Atoms)$-action that $\mathsf{FM}$
requires.

\subsection{Finite support is automatic}
\label{sec:fmmodel:support}

Write $\Red$ for the class of well-formed red sets, taken modulo red equality,
all of whose atoms lie in $\Atoms$. The finitary rules of the formal theory ---
$\mathrm{embrace}_i$ has finite arity, $\mathrm{Atomicity}$ adjoins one atom, and
$\rcup,\rcap,\rsm$ and $\mathrm{Comprehension}$ over finite sets preserve
finiteness --- yield only \emph{hereditarily finite} red sets. For $S\in\Red$ let
$\atomsof(S)\subseteq\Atoms$ be the finite set of atoms occurring hereditarily in
$S$.

\begin{lemma}[Support]
\label{lem:support}
For every $S\in\Red$ the finite set $\atomsof(S)$ supports $S$: if
$\pi\in\Sym(\Atoms)$ fixes $\atomsof(S)$ pointwise then $\pi\cdot S = S$.
Moreover $\atomsof(S)$ is the \emph{least} support, so $\supp(S)=\atomsof(S)$.
\end{lemma}

\begin{proof}[Proof sketch]
That $\atomsof(S)$ supports $S$ is an induction on the rank of $S$: an atom
$a\in\atomsof(S)$ is fixed by hypothesis, and for a set $S=\rbr{e_1,\dots,e_m}$ a
permutation fixing $\atomsof(S)\supseteq\bigcup_i \atomsof(e_i)$ fixes each $e_i$
by induction, hence fixes $S$. Leastness uses that $\Atoms$ is infinite: if
$a\in\atomsof(S)$ then, choosing a fresh $b\notin\atomsof(S)$, the transposition
$(a\,b)$ fixes $\atomsof(S)\setminus\{a\}$ yet moves $S$, so $a$ lies in every
support. Hence no proper subset of $\atomsof(S)$ supports $S$.
\end{proof}

\begin{corollary}
\label{cor:hfs}
Every $S\in\Red$ is hereditarily finitely supported, with
$\supp(S)=\atomsof(S)$. Freshness is occurrence-failure,
$a\mathrel{\#}S \iff a\notin\atomsof(S)$, and since $\Atoms$ is infinite the
fresh quantifier $\new$ and its some/any law (Section~\ref{sec:fm}) hold in
$\Red$.
\end{corollary}

The point worth stressing is that in $\mathsf{FM}$ finite support is an
\emph{axiom}, imposed by cutting the cumulative hierarchy down to
$\mathcal{U}_{\mathsf{FM}}$. Here it is a \emph{theorem}: the finitary red rules
never build anything that is not finitely supported, so the cut is vacuous and
nothing need be assumed.

\begin{remark}[Two points of precision]
\label{rem:group}
Lemma~\ref{lem:support} rests on two facts worth making explicit. First, the
group in play is the \emph{finitary} symmetric group of $\Atoms$ --- the
permutations generated by the transpositions that $\mathrm{swap}$ realises --- for
which supports are closed under intersection and least supports therefore exist;
on an arbitrary subgroup they need not. Second, the action permutes atoms while
treating each as structureless: it is an automorphism of the \emph{red} structure
$(\Red,\rin)$ alone, and is not asked to respect the black membership $\bin$
internal to an atom. The atoms are accordingly black-distinct but red-opaque, and
the entire support theory lives at the red level.
\end{remark}

\subsection{The representation theorem}
\label{sec:fmmodel:thm}

Let $\mathcal{M}$ be the structure with universe $\Atoms\cup\Red$, with the atoms
$\Atoms$ as urelements, with $\in^{\mathcal M}$ the red membership relation and
$=^{\mathcal M}$ red equality.

\begin{theorem}[Representation]
\label{thm:rep}
$\mathcal{M}$ is a model of hereditarily-finite Fraenkel--Mostowski set theory:
the $\mathsf{FM}$ analogue of $\mathrm{HF}(\Atoms)$. Concretely, interpreting the
atoms of $\mathsf{FM}$ as $\Atoms$, its sets as $\Red$, and its membership as red
$\rin$, the structure $\mathcal{M}$ validates extensionality, the empty set,
pairing, (finite) union, (finite) separation and foundation in their
$\mathsf{ZFA}$ form, carries the $\Sym(\Atoms)$-action of
Section~\ref{sec:fmmodel:action}, and satisfies the finite-support axiom. The
fresh quantifier and the some/any theorem hold throughout.
\end{theorem}

\begin{proof}[Proof sketch]
The atoms are membersless: red $\rin$ cannot enter a black set, so no
$a\in\Atoms$ has red elements, and atoms are $\in$-minimal and disjoint from
$\Red$ in sort, as $\mathsf{ZFA}$ demands. Foundation holds because each red set
is built at a finite rank over the atoms, so $\in^{\mathcal M}$ is well-founded.
The remaining axioms come straight from the rules of the formal theory:
$\mathrm{Foundation}$ gives $\remp$; $\mathrm{embrace}$ with $\mathrm{Union}$
gives pairing $\rbr{e_1}\rcup\rbr{e_2}$ and, iterated, every finite set;
$\mathrm{Union}$ gives binary and hence finite unions; $\mathrm{Subtraction}$ and
$\mathrm{Comprehension}$ over finite domains give finite separation; and
extensionality holds for the model provided provable red equality coincides with
sameness of members --- a property the equations of the formal theory are intended
to force, and which Section~\ref{sec:rigor} records as one of the obligations a
fully rigorous development must discharge. The $\Sym(\Atoms)$-action and the
finite-support axiom are Section~\ref{sec:fmmodel:action} and
Corollary~\ref{cor:hfs}; the some/any law follows because $\Atoms$ is infinite
and every $S\in\Red$ has finite support. Powerset and infinity are absent --- as
expected of an $\mathrm{HF}$ universe --- and are the subject of
Section~\ref{sec:fmmodel:infinitary}.
\end{proof}

\begin{corollary}[Two interlocked models]
\label{cor:two}
By the red/black symmetry of the construction, the same argument run with the
colours exchanged produces a black model $\bk{\mathcal M}$ of
hereditarily-finite $\mathsf{FM}$ whose atoms are the red numerals
$\{\,\rd{n} : n\in\omega\,\}\subseteq\rSetX$. The two models interlock: the atoms
of each are sets of the other. In the game-semantic reading of the
\emph{Intuitions} above, player's atoms are opponent's strategies and conversely
--- a single reflective universe carrying both $\mathsf{FM}$ models at once, each
underwriting the other's names.
\end{corollary}

\begin{remark}[What the $\mathrm{HF}$ restriction costs]
\label{rem:hftrivial}
It should be said plainly what is and is not delivered. The phenomena that give
$\mathsf{FM}$ its character --- finite support as a genuine constraint on
\emph{infinite} objects, equivariance with non-trivial orbits, and the some/any
law quantifying over infinitely many fresh names against a fixed infinite
support --- all require infinite, finitely-supported sets, of which $\mathcal{M}$
has none. On the hereditarily finite fragment these properties hold, but
trivially: every set is finite, so finite support is automatic and the some/any
law degenerates. $\mathcal{M}$ is thus best described as a model of the
\emph{nominal hereditarily finite sets}: faithful to the $\mathsf{ZFA}$ and
$\mathsf{FM}$ axioms, but exercising the distinctively nominal content only once
the infinitary extension of Section~\ref{sec:fmmodel:infinitary} is in force. We
retain the phrase ``model of $\mathsf{FM}$'' only under this qualification.
\end{remark}

\subsection{Finitary by design: the hereditarily finite core}
\label{sec:fmmodel:finitary}

It is worth being explicit about why Theorem~\ref{thm:rep} lands on the
hereditarily finite fragment, because the reason is also the source of the
construction's advantage. The signature of the formal theory offers, as set
constructors, only $\remp$, the finite-arity embraces $\mathrm{embrace}_i$
($|\mathrm{embrace}_i| = i < \omega$), the binary $\rcup,\rcap,\rsm$, and
$\mathrm{Comprehension}$ over finite domains $S_1\times\dots\times S_n$. Each of
these maps finite sets to finite sets, and the only base case is $\remp$;
so a trivial induction shows that \emph{every} derivable red set is hereditarily
finite. The model is $\mathrm{HF}(\Atoms)$, the nominal hereditarily finite sets
over the atoms.

Two of $\mathsf{FM}$'s defining commitments are, in this regime, theorems rather
than axioms.
\begin{itemize}
\item \textbf{Foundation} holds because every red set sits at a finite rank over
  the atoms, so $\in^{\mathcal{M}}$ is well-founded (Theorem~\ref{thm:rep}); we
  never assumed regularity.
\item \textbf{Finite support} holds because a hereditarily finite set mentions
  only finitely many atoms, each of which is its own least support
  (Lemma~\ref{lem:support}); we never assumed the finite-support axiom.
\end{itemize}
Correspondingly there is no infinity axiom and no infinite risk: the stream
$S[X]$ is read only through $\mathrm{take}(\Sigma,n)$ for finite $n$, so at no
point is an actually-infinite set formed. This is the regime in which the
reflective theory strictly improves on axiomatic $\mathsf{FM}$ --- both of the
moves that define $\mathsf{FM}$, foundation and the finite-support cut, are
\emph{derived} from the construction rather than \emph{posited}.

\subsection{Going infinitary, and the return of the cut}
\label{sec:fmmodel:infinitary}

Full $\mathsf{FM}$ also has infinite, finitely-supported sets --- $\Atoms$
itself, or the equivariant powerset $\pow_{\fs}(\Atoms)$ of finitely-supported
sets of atoms --- which the finitary signature cannot build. To reach them one
adds an \emph{infinitary} constructor. The most economical is an $\omega$-ary (or
$\kappa$-ary) embrace that consumes a whole stream rather than a finite prefix,
\[
  \inferrule{\Sigma;\Gamma\vdash S \\ (\sigma,\Sigma') = \mathrm{take}(\Sigma,\omega)}
            {\Sigma';\,\sigma,\Gamma\vdash \mathrm{embrace}_\omega(\sigma)\rcup S}
  \quad\textsc{Infinitary-Embrace},
\]
where $\mathrm{take}(\Sigma,\omega)$ returns the entire stream; equivalently one
adjoins a powerset operator and an infinity axiom and climbs the red cumulative
hierarchy $\mathcal{U}_\alpha(\Atoms)$ through all ordinals. Either way two things
change, and --- this is the point the finitary core conceals --- they change
\emph{together}.

\paragraph{An actual infinity, hence one $\omega$ of risk.}
The lazy stream becomes a completed set; $\Atoms$ itself becomes a red set. By
Remark~\ref{rem:risk} the commitment incurred is exactly black infinity, a single
$\omega$-style axiom. No primitive atoms are reintroduced --- the atoms remain
black hereditarily finite sets built from $\bemp$ --- but the potential infinity
of the finitary core is now actual.

\paragraph{Finite support is no longer free; the cut returns.}
Below $\omega$, finite support was automatic. The infinitary hierarchy is not so
forgiving: it contains sets that have \emph{no} finite support. The canonical
witness is the graph of the stream enumeration,
\[
  E \;=\; \{\,(n,a_n) : n\in\omega\,\},
  \qquad a_n = \text{the $n$-th atom of } S[X],
\]
which is fixed by a permutation $\pi$ only when $\pi$ preserves the enumeration,
i.e. only when $\pi$ is the identity on the infinitely many $a_n$; so no finite
set of atoms supports $E$. This is no accident. The formal theory already tells
us that $S[X]$ ``may be thought of as an ordering on $\Atoms$'', and a total order
on an infinite set of atoms is the textbook example of non-finitely-supported
data. The finitary rules never expose that order --- $\mathrm{take}$ reads only a
finite prefix --- which is precisely why support comes for free below $\omega$.
An infinitary embrace that swallows the whole stream imports the order, and with
it the non-finitely-supported sets.

Consequently, in the infinitary regime one must do exactly what axiomatic
$\mathsf{FM}$ does: cut down to the hereditarily finitely-supported part,
\[
  \mathcal{U}_{\mathsf{FM}} \;=\; \{\, x : x \text{ is hereditarily finitely supported}\,\},
\]
discarding $E$ and its kind. On the cut, the some/any theorem and the
equivariance principle hold verbatim, and the result satisfies the Gabbay--Pitts
axioms. The instructive point is that the cut is now the \emph{only} thing
borrowed from axiomatic $\mathsf{FM}$: the atoms still come from the dual copy,
the $\Sym(\Atoms)$-action still comes from $\mathrm{swap}$, and foundation still
comes from the construction (Section~\ref{sec:nwf}); only the finite-support
axiom must be re-imposed, and only because we elected to admit actual infinities.

This locates precisely where reflection saves risk and where it cannot. Atoms:
always free, from the dual copy. Foundation: always free, from local
well-foundedness (Section~\ref{sec:nwf}). Finite support: free in the finitary
core, an imposed cut once infinities are admitted. The finitary/infinitary
boundary is therefore not a technicality but the exact ledger line between what
the construction \emph{derives} and what it must still \emph{assume} --- and a
deliberate design choice for any theory built this way. We would keep the
finitary signature as the canonical core, where the accounting is cleanest, and
treat \textsc{Infinitary-Embrace} as an optional, clearly-priced extension.
).
%% rset.colors.tex
%% Red/black colour layer for the two intertwined copies, plus the FM symbols
%% used by rset.mainthm.tex and rset.applications.tex.
%%
%% EVERYTHING routes through \rd (red copy = player) and \bk (black copy =
%% opponent): redefine just those two macros to change the whole scheme (e.g.
%% to make the "black" copy dark blue for legibility).  Every definition is
%% \providecommand-guarded, so it is safe to \input this file from several
%% places and harmless if rset.local already defines these names.
%%
%% Requires \usepackage{xcolor} (already loaded in rset.tex).

% --- the two copies -------------------------------------------------------
\providecommand{\rd}[1]{\textcolor{red}{#1}}     % red copy  (player)
\providecommand{\bk}[1]{\textcolor{black}{#1}}   % black copy (opponent)

% --- theory names and atom supplies --------------------------------------
\providecommand{\rSetX}{\rd{\mathrm{Set}[X]}}    % red   Set[X]
\providecommand{\bSetX}{\bk{\mathrm{Set}[X]}}    % black Set[X]
\providecommand{\rX}{\rd{X}}                      % red   atom supply
\providecommand{\bX}{\bk{X}}                      % black atom supply

% --- empty set ------------------------------------------------------------
\providecommand{\remp}{\rd{\emptyset}}
\providecommand{\bemp}{\bk{\emptyset}}

% --- membership (relations) ----------------------------------------------
\providecommand{\rin}{\mathrel{\rd{\in}}}
\providecommand{\bin}{\mathrel{\bk{\in}}}
\providecommand{\rnin}{\mathrel{\rd{\notin}}}
\providecommand{\bnin}{\mathrel{\bk{\notin}}}

% --- binary operations ----------------------------------------------------
\providecommand{\rcup}{\mathbin{\rd{\cup}}}
\providecommand{\bcup}{\mathbin{\bk{\cup}}}
\providecommand{\rcap}{\mathbin{\rd{\cap}}}
\providecommand{\bcap}{\mathbin{\bk{\cap}}}
\providecommand{\rsm}{\mathbin{\rd{\setminus}}}
\providecommand{\bsm}{\mathbin{\bk{\setminus}}}

% --- comprehension braces of each theory ---------------------------------
\providecommand{\rbr}[1]{\rd{\{}\,#1\,\rd{\}}}
\providecommand{\bbr}[1]{\bk{\{}\,#1\,\bk{\}}}

% --- colour-blind membership (sees through the colour barrier) ------------
\providecommand{\cbin}{\mathrel{\in_{\!\star}}}

% --- FM symbols (shared with rset.stdfmset.tex) --------------------------
\providecommand{\Atoms}{\mathbb{A}}
\providecommand{\Sym}{\mathrm{Sym}}
\providecommand{\pow}{\mathcal{P}}
\providecommand{\supp}{\mathrm{supp}}
\providecommand{\fs}{\mathrm{fs}}
\providecommand{\new}{\mathbin{\reflectbox{\ensuremath{\mathsf{N}}}}}
\providecommand{\atomsof}{\mathrm{atoms}}
\providecommand{\Red}{\rd{\mathsf{R}}}           % the red HF universe
   % red/black colour layer + FM symbols (providecommand-guarded)

\section{Applications: nominal techniques founded on reflection}
\label{sec:applications}

The original motivation for $\mathsf{FM}$ in computer science was Pitts and
Gabbay's \emph{nominal} treatment of names and binding
\cite{DBLP:journals/fac/GabbayP02}. That treatment is founded squarely on
$\mathsf{FM}$: names are the atoms, $\alpha$-equivalence is an orbit relation
under the permutation action, and the freshness machinery is exactly support and
the $\new$-quantifier. Since Section~\ref{sec:fmmodel} realises $\mathsf{FM}$
inside the reflective universe, the entire nominal edifice descends with it. The
names that nominal techniques take as primitive become, here, black sets built
from the empty set; the reflective construction \emph{founds} the nominal theory
rather than presupposing its atoms.

\subsection{Nominal sets, internally}
\label{sec:applications:nom}

A \emph{nominal set} is a set carrying a $\Sym(\Atoms)$-action in which every
element is finitely supported; nominal sets and equivariant maps form the
category $\mathbf{Nom}$, equivalent to the Schanuel topos. The model
$\mathcal{M}$ of Theorem~\ref{thm:rep} is, on the nose, a nominal set: the action
is that of Section~\ref{sec:fmmodel:action}, and finite support is
Corollary~\ref{cor:hfs}. Every construction nominal practitioners rely on ---
the equivariance principle (anything definable from equivariant data is
equivariant), least supports, and the some/any law for $\new$ --- is therefore
available inside the reflective universe, with the colour discipline of
Section~\ref{sec:fmmodel} keeping the bookkeeping honest: names live in the
{\bk black} copy, the sets that bind them in the {\rd red} copy.

\subsection{Abstract syntax with binding}
\label{sec:applications:syntax}

The headline application of \cite{DBLP:journals/fac/GabbayP02} is a uniform
account of inductively-defined syntax \emph{modulo} $\alpha$-equivalence. Its key
device is \emph{name abstraction}: for a nominal set $X$,
\[
  [\Atoms]X \;=\; (\Atoms\times X)\big/\!\sim,
  \qquad
  \langle a\rangle x \sim \langle b\rangle y
  \;\;\Longleftrightarrow\;\;
  \new c.\ (a\,c)\cdot x = (b\,c)\cdot y,
\]
where $\langle a\rangle x$ is the class of $(a,x)$ and $(a\,c)$ is a
transposition. Both the action and $\new$ are present in $\mathcal{M}$, so
$[\Atoms]X$ is definable there, and it has the property that makes it a model of
binding: $\supp(\langle a\rangle x) = \supp(x)\setminus\{a\}$, i.e. abstraction
\emph{discharges} the bound name from the support, precisely as a binder removes
a variable from the free-variable set.

Inductive syntax is then obtained as usual. Object-level $\lambda$-terms, for
instance, are the least solution of
\[
  \Lambda \;\cong\; \Atoms \;+\; \Lambda\times\Lambda \;+\; [\Atoms]\Lambda
  \qquad(\text{variable, application, abstraction}),
\]
with $\alpha$-equivalence built into the third summand. Structural recursion and
induction come with the \emph{freshness condition for binders}, justified by the
some/any law: a definition or proof that treats a bound name as fresh for its
parameters is sound because, $\Atoms$ being infinite, ``some fresh name'' and
``every fresh name'' agree (Corollary~\ref{cor:hfs}). This is exactly the
reasoning principle nominal users invoke when they ``choose a fresh variable''.

The risk accounting carries through and sharpens the point. Each individual
term --- each element of $\Lambda$ --- is a hereditarily finite nominal set, so
it already lives in the $\mathrm{HF}$ model $\mathcal{M}$ of
Theorem~\ref{thm:rep}, founded on the empty set alone: a binding tree of any
fixed shape costs no infinite risk. The completed datatype $\Lambda$, an
infinite equivariant set, is one of the finitely-supported infinities of
Section~\ref{sec:fmmodel:infinitary}, and so costs exactly one $\omega$ --- black
infinity --- and still not a single primitive name. Where Pitts and Gabbay begin
with an infinite set of atoms given by fiat, here the names are manufactured, the
$\alpha$-structure is an orbit of the $\mathrm{swap}$ action, and the whole
apparatus is refinanced down to one empty set and one successor.

\subsection{Nominal algebraic theories}
\label{sec:applications:alg}

Beyond syntax, equational reasoning in the presence of names --- nominal
equational logic and the nominal Lawvere theories of Clouston
\cite{DBLP:journals/jcss/Clouston14} --- takes its semantics in $\mathbf{Nom}$.
As $\mathcal{M}$ furnishes $\mathbf{Nom}$-structure internally, any such theory
(name-restriction, scope extrusion, binding signatures and their equational
laws) can be interpreted in the reflective universe. The reflective setting adds
one structural feature these accounts do not have: by
Corollary~\ref{cor:two} there are \emph{two} interlocking name-universes, the red
and the black, each supplying the other's names. Reading red as player and black
as opponent recovers the game-semantic picture of the \emph{Intuitions}: a name
on one side is a strategy on the other, and a nominal theory and its dual coexist
in a single reflective universe.

\subsection{Summary}
\label{sec:applications:summary}

The chain is short and worth stating plainly. Nominal techniques rest on
$\mathsf{FM}$; $\mathsf{FM}$ embeds in the reflective universe
(Section~\ref{sec:fmmodel}); the reflective universe is built from the empty set.
Composing, the Pitts--Gabbay theory of names and binding is founded on two
recursively intertwined copies of ordinary set theory, with its primitive atoms
replaced by constructed black sets and its infinite risk replaced by the finite
risk of $\varnothing$ --- exactly the reconciliation announced in the
introduction, now delivered for the application that made $\mathsf{FM}$ matter to
computer science in the first place.
 (add the hook to rset.tex after
%% %% rset.mainthm.tex
%% A model of FM set theory inside the reflective theory of Sections 2--3,
%% with the two copies colour-coded (red = player, black = opponent).
%% Drop in via %% rset.mainthm.tex
%% A model of FM set theory inside the reflective theory of Sections 2--3,
%% with the two copies colour-coded (red = player, black = opponent).
%% Drop in via \input{rset.mainthm} (the hook already exists in rset.tex).
\input{rset.colors}   % red/black colour layer + FM symbols (providecommand-guarded)

\section{A model of $\mathsf{FM}$ in the reflective universe}
\label{sec:fmmodel}

We can now collect the dividend of the construction \emph{Tying the second
recursive knot} above. The reflective theory was assembled from the empty set
alone, yet --- we claim --- it already contains a model of $\mathsf{FM}$, and in
fact two of them, interlocked. The atoms such a model needs are not postulated;
they are \emph{supplied by the dual copy}. Throughout we fix the two tied
theories introduced above, the red copy $\rSetX$ and the black copy $\bSetX$,
with
\[
  \rX \;=\; \bSetX, \qquad \bX \;=\; \rSetX,
\]
so that the admissible atoms of the red theory are exactly the black sets, and
\emph{vice versa}. We write the red copy in {\rd red} and the black copy in
{\bk black} throughout, so that a red comprehension $\rbr{\cdots}$ and a black
comprehension $\bbr{\cdots}$, or the red and black membership relations
$\rin$ and $\bin$, are told apart at a glance.

\subsection{Atoms out of the empty set}
\label{sec:fmmodel:atoms}

Inside the black copy, build the von Neumann numerals from the black empty set,
\[
  \bk{0} \;=\; \bemp,
  \qquad
  \bk{n+1} \;=\; \bk{n} \bcup \bbr{\bk{n}},
\]
each a black, hereditarily finite set obtained, ultimately, from $\bemp$. By
black extensionality (the equations of the formal theory above) these are
pairwise distinct: $\bk{n} \neq \bk{m}$ whenever $n\neq m$. Put
\[
  \Atoms \;:=\; \{\, \bk{n} \;:\; n\in\omega \,\}\ \subseteq\ \bSetX .
\]

Two observations make $\Atoms$ behave \emph{as a set of atoms} for the red
theory.
\begin{itemize}
\item \textbf{Admissibility.} By the second recursive knot $\rX=\bSetX$, every
  black set is an admissible red atom; in particular each $\bk{n}$ is one.
\item \textbf{Opacity.} The red membership relation $\rin$ cannot see into a
  black set (the \emph{Intuitions} above). Hence, \emph{as far as the red theory
  can tell}, an element $\bk{n}$ of $\Atoms$ has no internal structure
  whatsoever: the only facts the red theory may assert about $\bk{n}$ and
  $\bk{m}$ are $\bk{n} = \bk{m}$ or $\bk{n} \neq \bk{m}$, and those are settled,
  out of sight, by black equality.
\end{itemize}
Thus $\Atoms$ presents to the red theory precisely as $\mathsf{FM}$'s set of
atoms does: a countable family of mutually distinct, structureless objects that
may sit inside sets but admit no membership query of their own.

\begin{remark}[Where the risk went]
\label{rem:risk}
This is the sense, promised in the introduction, in which atoms cost no infinite
risk. $\mathsf{ZFA}$ and $\mathsf{FM}$ \emph{postulate} an infinite set of
primitive atoms; the postulate is the risk. Here nothing is postulated beyond
$\bemp$. The atoms are black hereditarily finite sets, and the stream $S[X]$ of
the formal theory hands them out lazily, $\mathrm{take}(\Sigma,n)$ at a time, so
that any red set ever built mentions only finitely many of them. The supply is
\emph{potentially} infinite, never an actual completed infinity: the commitment
is the black empty set together with the same successor step that underwrites
$\omega$, and not one primitive atom more.
\end{remark}

\subsection{The permutation action}
\label{sec:fmmodel:action}

The stream operation $\mathrm{swap}(S[X],\rho)$ of the formal theory applies a
finite permutation $\rho=\{\,\bk{j_1}/\bk{i_1}:\dots:\bk{j_k}/\bk{i_k}\,\}$ to
the ordering of atoms; these finite permutations generate the full group
$\Sym(\Atoms)$. A permutation $\pi\in\Sym(\Atoms)$ acts on red sets by
red-membership recursion exactly as in $\mathsf{FM}$,
\[
  \pi\cdot a = \pi(a)\ (a\in\Atoms),
  \qquad
  \pi\cdot S = \rd{\{}\, \pi\cdot e \;:\; e \rin S \,\rd{\}}.
\]
Because $\pi$ moves only atoms and never the black interior of an atom, and
because the red operations $\remp,\ \rcup,\ \rcap,\ \rsm$ and
$\mathrm{embrace}_i$ are visibly equivariant (a renaming commutes with each), the
action is by automorphisms of the red theory and respects the equations of the
formal theory. We have, verbatim, the $\Sym(\Atoms)$-action that $\mathsf{FM}$
requires.

\subsection{Finite support is automatic}
\label{sec:fmmodel:support}

Write $\Red$ for the class of well-formed red sets, taken modulo red equality,
all of whose atoms lie in $\Atoms$. The finitary rules of the formal theory ---
$\mathrm{embrace}_i$ has finite arity, $\mathrm{Atomicity}$ adjoins one atom, and
$\rcup,\rcap,\rsm$ and $\mathrm{Comprehension}$ over finite sets preserve
finiteness --- yield only \emph{hereditarily finite} red sets. For $S\in\Red$ let
$\atomsof(S)\subseteq\Atoms$ be the finite set of atoms occurring hereditarily in
$S$.

\begin{lemma}[Support]
\label{lem:support}
For every $S\in\Red$ the finite set $\atomsof(S)$ supports $S$: if
$\pi\in\Sym(\Atoms)$ fixes $\atomsof(S)$ pointwise then $\pi\cdot S = S$.
Moreover $\atomsof(S)$ is the \emph{least} support, so $\supp(S)=\atomsof(S)$.
\end{lemma}

\begin{proof}[Proof sketch]
That $\atomsof(S)$ supports $S$ is an induction on the rank of $S$: an atom
$a\in\atomsof(S)$ is fixed by hypothesis, and for a set $S=\rbr{e_1,\dots,e_m}$ a
permutation fixing $\atomsof(S)\supseteq\bigcup_i \atomsof(e_i)$ fixes each $e_i$
by induction, hence fixes $S$. Leastness uses that $\Atoms$ is infinite: if
$a\in\atomsof(S)$ then, choosing a fresh $b\notin\atomsof(S)$, the transposition
$(a\,b)$ fixes $\atomsof(S)\setminus\{a\}$ yet moves $S$, so $a$ lies in every
support. Hence no proper subset of $\atomsof(S)$ supports $S$.
\end{proof}

\begin{corollary}
\label{cor:hfs}
Every $S\in\Red$ is hereditarily finitely supported, with
$\supp(S)=\atomsof(S)$. Freshness is occurrence-failure,
$a\mathrel{\#}S \iff a\notin\atomsof(S)$, and since $\Atoms$ is infinite the
fresh quantifier $\new$ and its some/any law (Section~\ref{sec:fm}) hold in
$\Red$.
\end{corollary}

The point worth stressing is that in $\mathsf{FM}$ finite support is an
\emph{axiom}, imposed by cutting the cumulative hierarchy down to
$\mathcal{U}_{\mathsf{FM}}$. Here it is a \emph{theorem}: the finitary red rules
never build anything that is not finitely supported, so the cut is vacuous and
nothing need be assumed.

\begin{remark}[Two points of precision]
\label{rem:group}
Lemma~\ref{lem:support} rests on two facts worth making explicit. First, the
group in play is the \emph{finitary} symmetric group of $\Atoms$ --- the
permutations generated by the transpositions that $\mathrm{swap}$ realises --- for
which supports are closed under intersection and least supports therefore exist;
on an arbitrary subgroup they need not. Second, the action permutes atoms while
treating each as structureless: it is an automorphism of the \emph{red} structure
$(\Red,\rin)$ alone, and is not asked to respect the black membership $\bin$
internal to an atom. The atoms are accordingly black-distinct but red-opaque, and
the entire support theory lives at the red level.
\end{remark}

\subsection{The representation theorem}
\label{sec:fmmodel:thm}

Let $\mathcal{M}$ be the structure with universe $\Atoms\cup\Red$, with the atoms
$\Atoms$ as urelements, with $\in^{\mathcal M}$ the red membership relation and
$=^{\mathcal M}$ red equality.

\begin{theorem}[Representation]
\label{thm:rep}
$\mathcal{M}$ is a model of hereditarily-finite Fraenkel--Mostowski set theory:
the $\mathsf{FM}$ analogue of $\mathrm{HF}(\Atoms)$. Concretely, interpreting the
atoms of $\mathsf{FM}$ as $\Atoms$, its sets as $\Red$, and its membership as red
$\rin$, the structure $\mathcal{M}$ validates extensionality, the empty set,
pairing, (finite) union, (finite) separation and foundation in their
$\mathsf{ZFA}$ form, carries the $\Sym(\Atoms)$-action of
Section~\ref{sec:fmmodel:action}, and satisfies the finite-support axiom. The
fresh quantifier and the some/any theorem hold throughout.
\end{theorem}

\begin{proof}[Proof sketch]
The atoms are membersless: red $\rin$ cannot enter a black set, so no
$a\in\Atoms$ has red elements, and atoms are $\in$-minimal and disjoint from
$\Red$ in sort, as $\mathsf{ZFA}$ demands. Foundation holds because each red set
is built at a finite rank over the atoms, so $\in^{\mathcal M}$ is well-founded.
The remaining axioms come straight from the rules of the formal theory:
$\mathrm{Foundation}$ gives $\remp$; $\mathrm{embrace}$ with $\mathrm{Union}$
gives pairing $\rbr{e_1}\rcup\rbr{e_2}$ and, iterated, every finite set;
$\mathrm{Union}$ gives binary and hence finite unions; $\mathrm{Subtraction}$ and
$\mathrm{Comprehension}$ over finite domains give finite separation; and
extensionality holds for the model provided provable red equality coincides with
sameness of members --- a property the equations of the formal theory are intended
to force, and which Section~\ref{sec:rigor} records as one of the obligations a
fully rigorous development must discharge. The $\Sym(\Atoms)$-action and the
finite-support axiom are Section~\ref{sec:fmmodel:action} and
Corollary~\ref{cor:hfs}; the some/any law follows because $\Atoms$ is infinite
and every $S\in\Red$ has finite support. Powerset and infinity are absent --- as
expected of an $\mathrm{HF}$ universe --- and are the subject of
Section~\ref{sec:fmmodel:infinitary}.
\end{proof}

\begin{corollary}[Two interlocked models]
\label{cor:two}
By the red/black symmetry of the construction, the same argument run with the
colours exchanged produces a black model $\bk{\mathcal M}$ of
hereditarily-finite $\mathsf{FM}$ whose atoms are the red numerals
$\{\,\rd{n} : n\in\omega\,\}\subseteq\rSetX$. The two models interlock: the atoms
of each are sets of the other. In the game-semantic reading of the
\emph{Intuitions} above, player's atoms are opponent's strategies and conversely
--- a single reflective universe carrying both $\mathsf{FM}$ models at once, each
underwriting the other's names.
\end{corollary}

\begin{remark}[What the $\mathrm{HF}$ restriction costs]
\label{rem:hftrivial}
It should be said plainly what is and is not delivered. The phenomena that give
$\mathsf{FM}$ its character --- finite support as a genuine constraint on
\emph{infinite} objects, equivariance with non-trivial orbits, and the some/any
law quantifying over infinitely many fresh names against a fixed infinite
support --- all require infinite, finitely-supported sets, of which $\mathcal{M}$
has none. On the hereditarily finite fragment these properties hold, but
trivially: every set is finite, so finite support is automatic and the some/any
law degenerates. $\mathcal{M}$ is thus best described as a model of the
\emph{nominal hereditarily finite sets}: faithful to the $\mathsf{ZFA}$ and
$\mathsf{FM}$ axioms, but exercising the distinctively nominal content only once
the infinitary extension of Section~\ref{sec:fmmodel:infinitary} is in force. We
retain the phrase ``model of $\mathsf{FM}$'' only under this qualification.
\end{remark}

\subsection{Finitary by design: the hereditarily finite core}
\label{sec:fmmodel:finitary}

It is worth being explicit about why Theorem~\ref{thm:rep} lands on the
hereditarily finite fragment, because the reason is also the source of the
construction's advantage. The signature of the formal theory offers, as set
constructors, only $\remp$, the finite-arity embraces $\mathrm{embrace}_i$
($|\mathrm{embrace}_i| = i < \omega$), the binary $\rcup,\rcap,\rsm$, and
$\mathrm{Comprehension}$ over finite domains $S_1\times\dots\times S_n$. Each of
these maps finite sets to finite sets, and the only base case is $\remp$;
so a trivial induction shows that \emph{every} derivable red set is hereditarily
finite. The model is $\mathrm{HF}(\Atoms)$, the nominal hereditarily finite sets
over the atoms.

Two of $\mathsf{FM}$'s defining commitments are, in this regime, theorems rather
than axioms.
\begin{itemize}
\item \textbf{Foundation} holds because every red set sits at a finite rank over
  the atoms, so $\in^{\mathcal{M}}$ is well-founded (Theorem~\ref{thm:rep}); we
  never assumed regularity.
\item \textbf{Finite support} holds because a hereditarily finite set mentions
  only finitely many atoms, each of which is its own least support
  (Lemma~\ref{lem:support}); we never assumed the finite-support axiom.
\end{itemize}
Correspondingly there is no infinity axiom and no infinite risk: the stream
$S[X]$ is read only through $\mathrm{take}(\Sigma,n)$ for finite $n$, so at no
point is an actually-infinite set formed. This is the regime in which the
reflective theory strictly improves on axiomatic $\mathsf{FM}$ --- both of the
moves that define $\mathsf{FM}$, foundation and the finite-support cut, are
\emph{derived} from the construction rather than \emph{posited}.

\subsection{Going infinitary, and the return of the cut}
\label{sec:fmmodel:infinitary}

Full $\mathsf{FM}$ also has infinite, finitely-supported sets --- $\Atoms$
itself, or the equivariant powerset $\pow_{\fs}(\Atoms)$ of finitely-supported
sets of atoms --- which the finitary signature cannot build. To reach them one
adds an \emph{infinitary} constructor. The most economical is an $\omega$-ary (or
$\kappa$-ary) embrace that consumes a whole stream rather than a finite prefix,
\[
  \inferrule{\Sigma;\Gamma\vdash S \\ (\sigma,\Sigma') = \mathrm{take}(\Sigma,\omega)}
            {\Sigma';\,\sigma,\Gamma\vdash \mathrm{embrace}_\omega(\sigma)\rcup S}
  \quad\textsc{Infinitary-Embrace},
\]
where $\mathrm{take}(\Sigma,\omega)$ returns the entire stream; equivalently one
adjoins a powerset operator and an infinity axiom and climbs the red cumulative
hierarchy $\mathcal{U}_\alpha(\Atoms)$ through all ordinals. Either way two things
change, and --- this is the point the finitary core conceals --- they change
\emph{together}.

\paragraph{An actual infinity, hence one $\omega$ of risk.}
The lazy stream becomes a completed set; $\Atoms$ itself becomes a red set. By
Remark~\ref{rem:risk} the commitment incurred is exactly black infinity, a single
$\omega$-style axiom. No primitive atoms are reintroduced --- the atoms remain
black hereditarily finite sets built from $\bemp$ --- but the potential infinity
of the finitary core is now actual.

\paragraph{Finite support is no longer free; the cut returns.}
Below $\omega$, finite support was automatic. The infinitary hierarchy is not so
forgiving: it contains sets that have \emph{no} finite support. The canonical
witness is the graph of the stream enumeration,
\[
  E \;=\; \{\,(n,a_n) : n\in\omega\,\},
  \qquad a_n = \text{the $n$-th atom of } S[X],
\]
which is fixed by a permutation $\pi$ only when $\pi$ preserves the enumeration,
i.e. only when $\pi$ is the identity on the infinitely many $a_n$; so no finite
set of atoms supports $E$. This is no accident. The formal theory already tells
us that $S[X]$ ``may be thought of as an ordering on $\Atoms$'', and a total order
on an infinite set of atoms is the textbook example of non-finitely-supported
data. The finitary rules never expose that order --- $\mathrm{take}$ reads only a
finite prefix --- which is precisely why support comes for free below $\omega$.
An infinitary embrace that swallows the whole stream imports the order, and with
it the non-finitely-supported sets.

Consequently, in the infinitary regime one must do exactly what axiomatic
$\mathsf{FM}$ does: cut down to the hereditarily finitely-supported part,
\[
  \mathcal{U}_{\mathsf{FM}} \;=\; \{\, x : x \text{ is hereditarily finitely supported}\,\},
\]
discarding $E$ and its kind. On the cut, the some/any theorem and the
equivariance principle hold verbatim, and the result satisfies the Gabbay--Pitts
axioms. The instructive point is that the cut is now the \emph{only} thing
borrowed from axiomatic $\mathsf{FM}$: the atoms still come from the dual copy,
the $\Sym(\Atoms)$-action still comes from $\mathrm{swap}$, and foundation still
comes from the construction (Section~\ref{sec:nwf}); only the finite-support
axiom must be re-imposed, and only because we elected to admit actual infinities.

This locates precisely where reflection saves risk and where it cannot. Atoms:
always free, from the dual copy. Foundation: always free, from local
well-foundedness (Section~\ref{sec:nwf}). Finite support: free in the finitary
core, an imposed cut once infinities are admitted. The finitary/infinitary
boundary is therefore not a technicality but the exact ledger line between what
the construction \emph{derives} and what it must still \emph{assume} --- and a
deliberate design choice for any theory built this way. We would keep the
finitary signature as the canonical core, where the accounting is cleanest, and
treat \textsc{Infinitary-Embrace} as an optional, clearly-priced extension.
 (the hook already exists in rset.tex).
%% rset.colors.tex
%% Red/black colour layer for the two intertwined copies, plus the FM symbols
%% used by rset.mainthm.tex and rset.applications.tex.
%%
%% EVERYTHING routes through \rd (red copy = player) and \bk (black copy =
%% opponent): redefine just those two macros to change the whole scheme (e.g.
%% to make the "black" copy dark blue for legibility).  Every definition is
%% \providecommand-guarded, so it is safe to \input this file from several
%% places and harmless if rset.local already defines these names.
%%
%% Requires \usepackage{xcolor} (already loaded in rset.tex).

% --- the two copies -------------------------------------------------------
\providecommand{\rd}[1]{\textcolor{red}{#1}}     % red copy  (player)
\providecommand{\bk}[1]{\textcolor{black}{#1}}   % black copy (opponent)

% --- theory names and atom supplies --------------------------------------
\providecommand{\rSetX}{\rd{\mathrm{Set}[X]}}    % red   Set[X]
\providecommand{\bSetX}{\bk{\mathrm{Set}[X]}}    % black Set[X]
\providecommand{\rX}{\rd{X}}                      % red   atom supply
\providecommand{\bX}{\bk{X}}                      % black atom supply

% --- empty set ------------------------------------------------------------
\providecommand{\remp}{\rd{\emptyset}}
\providecommand{\bemp}{\bk{\emptyset}}

% --- membership (relations) ----------------------------------------------
\providecommand{\rin}{\mathrel{\rd{\in}}}
\providecommand{\bin}{\mathrel{\bk{\in}}}
\providecommand{\rnin}{\mathrel{\rd{\notin}}}
\providecommand{\bnin}{\mathrel{\bk{\notin}}}

% --- binary operations ----------------------------------------------------
\providecommand{\rcup}{\mathbin{\rd{\cup}}}
\providecommand{\bcup}{\mathbin{\bk{\cup}}}
\providecommand{\rcap}{\mathbin{\rd{\cap}}}
\providecommand{\bcap}{\mathbin{\bk{\cap}}}
\providecommand{\rsm}{\mathbin{\rd{\setminus}}}
\providecommand{\bsm}{\mathbin{\bk{\setminus}}}

% --- comprehension braces of each theory ---------------------------------
\providecommand{\rbr}[1]{\rd{\{}\,#1\,\rd{\}}}
\providecommand{\bbr}[1]{\bk{\{}\,#1\,\bk{\}}}

% --- colour-blind membership (sees through the colour barrier) ------------
\providecommand{\cbin}{\mathrel{\in_{\!\star}}}

% --- FM symbols (shared with rset.stdfmset.tex) --------------------------
\providecommand{\Atoms}{\mathbb{A}}
\providecommand{\Sym}{\mathrm{Sym}}
\providecommand{\pow}{\mathcal{P}}
\providecommand{\supp}{\mathrm{supp}}
\providecommand{\fs}{\mathrm{fs}}
\providecommand{\new}{\mathbin{\reflectbox{\ensuremath{\mathsf{N}}}}}
\providecommand{\atomsof}{\mathrm{atoms}}
\providecommand{\Red}{\rd{\mathsf{R}}}           % the red HF universe
   % red/black colour layer + FM symbols (providecommand-guarded)

\section{A model of $\mathsf{FM}$ in the reflective universe}
\label{sec:fmmodel}

We can now collect the dividend of the construction \emph{Tying the second
recursive knot} above. The reflective theory was assembled from the empty set
alone, yet --- we claim --- it already contains a model of $\mathsf{FM}$, and in
fact two of them, interlocked. The atoms such a model needs are not postulated;
they are \emph{supplied by the dual copy}. Throughout we fix the two tied
theories introduced above, the red copy $\rSetX$ and the black copy $\bSetX$,
with
\[
  \rX \;=\; \bSetX, \qquad \bX \;=\; \rSetX,
\]
so that the admissible atoms of the red theory are exactly the black sets, and
\emph{vice versa}. We write the red copy in {\rd red} and the black copy in
{\bk black} throughout, so that a red comprehension $\rbr{\cdots}$ and a black
comprehension $\bbr{\cdots}$, or the red and black membership relations
$\rin$ and $\bin$, are told apart at a glance.

\subsection{Atoms out of the empty set}
\label{sec:fmmodel:atoms}

Inside the black copy, build the von Neumann numerals from the black empty set,
\[
  \bk{0} \;=\; \bemp,
  \qquad
  \bk{n+1} \;=\; \bk{n} \bcup \bbr{\bk{n}},
\]
each a black, hereditarily finite set obtained, ultimately, from $\bemp$. By
black extensionality (the equations of the formal theory above) these are
pairwise distinct: $\bk{n} \neq \bk{m}$ whenever $n\neq m$. Put
\[
  \Atoms \;:=\; \{\, \bk{n} \;:\; n\in\omega \,\}\ \subseteq\ \bSetX .
\]

Two observations make $\Atoms$ behave \emph{as a set of atoms} for the red
theory.
\begin{itemize}
\item \textbf{Admissibility.} By the second recursive knot $\rX=\bSetX$, every
  black set is an admissible red atom; in particular each $\bk{n}$ is one.
\item \textbf{Opacity.} The red membership relation $\rin$ cannot see into a
  black set (the \emph{Intuitions} above). Hence, \emph{as far as the red theory
  can tell}, an element $\bk{n}$ of $\Atoms$ has no internal structure
  whatsoever: the only facts the red theory may assert about $\bk{n}$ and
  $\bk{m}$ are $\bk{n} = \bk{m}$ or $\bk{n} \neq \bk{m}$, and those are settled,
  out of sight, by black equality.
\end{itemize}
Thus $\Atoms$ presents to the red theory precisely as $\mathsf{FM}$'s set of
atoms does: a countable family of mutually distinct, structureless objects that
may sit inside sets but admit no membership query of their own.

\begin{remark}[Where the risk went]
\label{rem:risk}
This is the sense, promised in the introduction, in which atoms cost no infinite
risk. $\mathsf{ZFA}$ and $\mathsf{FM}$ \emph{postulate} an infinite set of
primitive atoms; the postulate is the risk. Here nothing is postulated beyond
$\bemp$. The atoms are black hereditarily finite sets, and the stream $S[X]$ of
the formal theory hands them out lazily, $\mathrm{take}(\Sigma,n)$ at a time, so
that any red set ever built mentions only finitely many of them. The supply is
\emph{potentially} infinite, never an actual completed infinity: the commitment
is the black empty set together with the same successor step that underwrites
$\omega$, and not one primitive atom more.
\end{remark}

\subsection{The permutation action}
\label{sec:fmmodel:action}

The stream operation $\mathrm{swap}(S[X],\rho)$ of the formal theory applies a
finite permutation $\rho=\{\,\bk{j_1}/\bk{i_1}:\dots:\bk{j_k}/\bk{i_k}\,\}$ to
the ordering of atoms; these finite permutations generate the full group
$\Sym(\Atoms)$. A permutation $\pi\in\Sym(\Atoms)$ acts on red sets by
red-membership recursion exactly as in $\mathsf{FM}$,
\[
  \pi\cdot a = \pi(a)\ (a\in\Atoms),
  \qquad
  \pi\cdot S = \rd{\{}\, \pi\cdot e \;:\; e \rin S \,\rd{\}}.
\]
Because $\pi$ moves only atoms and never the black interior of an atom, and
because the red operations $\remp,\ \rcup,\ \rcap,\ \rsm$ and
$\mathrm{embrace}_i$ are visibly equivariant (a renaming commutes with each), the
action is by automorphisms of the red theory and respects the equations of the
formal theory. We have, verbatim, the $\Sym(\Atoms)$-action that $\mathsf{FM}$
requires.

\subsection{Finite support is automatic}
\label{sec:fmmodel:support}

Write $\Red$ for the class of well-formed red sets, taken modulo red equality,
all of whose atoms lie in $\Atoms$. The finitary rules of the formal theory ---
$\mathrm{embrace}_i$ has finite arity, $\mathrm{Atomicity}$ adjoins one atom, and
$\rcup,\rcap,\rsm$ and $\mathrm{Comprehension}$ over finite sets preserve
finiteness --- yield only \emph{hereditarily finite} red sets. For $S\in\Red$ let
$\atomsof(S)\subseteq\Atoms$ be the finite set of atoms occurring hereditarily in
$S$.

\begin{lemma}[Support]
\label{lem:support}
For every $S\in\Red$ the finite set $\atomsof(S)$ supports $S$: if
$\pi\in\Sym(\Atoms)$ fixes $\atomsof(S)$ pointwise then $\pi\cdot S = S$.
Moreover $\atomsof(S)$ is the \emph{least} support, so $\supp(S)=\atomsof(S)$.
\end{lemma}

\begin{proof}[Proof sketch]
That $\atomsof(S)$ supports $S$ is an induction on the rank of $S$: an atom
$a\in\atomsof(S)$ is fixed by hypothesis, and for a set $S=\rbr{e_1,\dots,e_m}$ a
permutation fixing $\atomsof(S)\supseteq\bigcup_i \atomsof(e_i)$ fixes each $e_i$
by induction, hence fixes $S$. Leastness uses that $\Atoms$ is infinite: if
$a\in\atomsof(S)$ then, choosing a fresh $b\notin\atomsof(S)$, the transposition
$(a\,b)$ fixes $\atomsof(S)\setminus\{a\}$ yet moves $S$, so $a$ lies in every
support. Hence no proper subset of $\atomsof(S)$ supports $S$.
\end{proof}

\begin{corollary}
\label{cor:hfs}
Every $S\in\Red$ is hereditarily finitely supported, with
$\supp(S)=\atomsof(S)$. Freshness is occurrence-failure,
$a\mathrel{\#}S \iff a\notin\atomsof(S)$, and since $\Atoms$ is infinite the
fresh quantifier $\new$ and its some/any law (Section~\ref{sec:fm}) hold in
$\Red$.
\end{corollary}

The point worth stressing is that in $\mathsf{FM}$ finite support is an
\emph{axiom}, imposed by cutting the cumulative hierarchy down to
$\mathcal{U}_{\mathsf{FM}}$. Here it is a \emph{theorem}: the finitary red rules
never build anything that is not finitely supported, so the cut is vacuous and
nothing need be assumed.

\begin{remark}[Two points of precision]
\label{rem:group}
Lemma~\ref{lem:support} rests on two facts worth making explicit. First, the
group in play is the \emph{finitary} symmetric group of $\Atoms$ --- the
permutations generated by the transpositions that $\mathrm{swap}$ realises --- for
which supports are closed under intersection and least supports therefore exist;
on an arbitrary subgroup they need not. Second, the action permutes atoms while
treating each as structureless: it is an automorphism of the \emph{red} structure
$(\Red,\rin)$ alone, and is not asked to respect the black membership $\bin$
internal to an atom. The atoms are accordingly black-distinct but red-opaque, and
the entire support theory lives at the red level.
\end{remark}

\subsection{The representation theorem}
\label{sec:fmmodel:thm}

Let $\mathcal{M}$ be the structure with universe $\Atoms\cup\Red$, with the atoms
$\Atoms$ as urelements, with $\in^{\mathcal M}$ the red membership relation and
$=^{\mathcal M}$ red equality.

\begin{theorem}[Representation]
\label{thm:rep}
$\mathcal{M}$ is a model of hereditarily-finite Fraenkel--Mostowski set theory:
the $\mathsf{FM}$ analogue of $\mathrm{HF}(\Atoms)$. Concretely, interpreting the
atoms of $\mathsf{FM}$ as $\Atoms$, its sets as $\Red$, and its membership as red
$\rin$, the structure $\mathcal{M}$ validates extensionality, the empty set,
pairing, (finite) union, (finite) separation and foundation in their
$\mathsf{ZFA}$ form, carries the $\Sym(\Atoms)$-action of
Section~\ref{sec:fmmodel:action}, and satisfies the finite-support axiom. The
fresh quantifier and the some/any theorem hold throughout.
\end{theorem}

\begin{proof}[Proof sketch]
The atoms are membersless: red $\rin$ cannot enter a black set, so no
$a\in\Atoms$ has red elements, and atoms are $\in$-minimal and disjoint from
$\Red$ in sort, as $\mathsf{ZFA}$ demands. Foundation holds because each red set
is built at a finite rank over the atoms, so $\in^{\mathcal M}$ is well-founded.
The remaining axioms come straight from the rules of the formal theory:
$\mathrm{Foundation}$ gives $\remp$; $\mathrm{embrace}$ with $\mathrm{Union}$
gives pairing $\rbr{e_1}\rcup\rbr{e_2}$ and, iterated, every finite set;
$\mathrm{Union}$ gives binary and hence finite unions; $\mathrm{Subtraction}$ and
$\mathrm{Comprehension}$ over finite domains give finite separation; and
extensionality holds for the model provided provable red equality coincides with
sameness of members --- a property the equations of the formal theory are intended
to force, and which Section~\ref{sec:rigor} records as one of the obligations a
fully rigorous development must discharge. The $\Sym(\Atoms)$-action and the
finite-support axiom are Section~\ref{sec:fmmodel:action} and
Corollary~\ref{cor:hfs}; the some/any law follows because $\Atoms$ is infinite
and every $S\in\Red$ has finite support. Powerset and infinity are absent --- as
expected of an $\mathrm{HF}$ universe --- and are the subject of
Section~\ref{sec:fmmodel:infinitary}.
\end{proof}

\begin{corollary}[Two interlocked models]
\label{cor:two}
By the red/black symmetry of the construction, the same argument run with the
colours exchanged produces a black model $\bk{\mathcal M}$ of
hereditarily-finite $\mathsf{FM}$ whose atoms are the red numerals
$\{\,\rd{n} : n\in\omega\,\}\subseteq\rSetX$. The two models interlock: the atoms
of each are sets of the other. In the game-semantic reading of the
\emph{Intuitions} above, player's atoms are opponent's strategies and conversely
--- a single reflective universe carrying both $\mathsf{FM}$ models at once, each
underwriting the other's names.
\end{corollary}

\begin{remark}[What the $\mathrm{HF}$ restriction costs]
\label{rem:hftrivial}
It should be said plainly what is and is not delivered. The phenomena that give
$\mathsf{FM}$ its character --- finite support as a genuine constraint on
\emph{infinite} objects, equivariance with non-trivial orbits, and the some/any
law quantifying over infinitely many fresh names against a fixed infinite
support --- all require infinite, finitely-supported sets, of which $\mathcal{M}$
has none. On the hereditarily finite fragment these properties hold, but
trivially: every set is finite, so finite support is automatic and the some/any
law degenerates. $\mathcal{M}$ is thus best described as a model of the
\emph{nominal hereditarily finite sets}: faithful to the $\mathsf{ZFA}$ and
$\mathsf{FM}$ axioms, but exercising the distinctively nominal content only once
the infinitary extension of Section~\ref{sec:fmmodel:infinitary} is in force. We
retain the phrase ``model of $\mathsf{FM}$'' only under this qualification.
\end{remark}

\subsection{Finitary by design: the hereditarily finite core}
\label{sec:fmmodel:finitary}

It is worth being explicit about why Theorem~\ref{thm:rep} lands on the
hereditarily finite fragment, because the reason is also the source of the
construction's advantage. The signature of the formal theory offers, as set
constructors, only $\remp$, the finite-arity embraces $\mathrm{embrace}_i$
($|\mathrm{embrace}_i| = i < \omega$), the binary $\rcup,\rcap,\rsm$, and
$\mathrm{Comprehension}$ over finite domains $S_1\times\dots\times S_n$. Each of
these maps finite sets to finite sets, and the only base case is $\remp$;
so a trivial induction shows that \emph{every} derivable red set is hereditarily
finite. The model is $\mathrm{HF}(\Atoms)$, the nominal hereditarily finite sets
over the atoms.

Two of $\mathsf{FM}$'s defining commitments are, in this regime, theorems rather
than axioms.
\begin{itemize}
\item \textbf{Foundation} holds because every red set sits at a finite rank over
  the atoms, so $\in^{\mathcal{M}}$ is well-founded (Theorem~\ref{thm:rep}); we
  never assumed regularity.
\item \textbf{Finite support} holds because a hereditarily finite set mentions
  only finitely many atoms, each of which is its own least support
  (Lemma~\ref{lem:support}); we never assumed the finite-support axiom.
\end{itemize}
Correspondingly there is no infinity axiom and no infinite risk: the stream
$S[X]$ is read only through $\mathrm{take}(\Sigma,n)$ for finite $n$, so at no
point is an actually-infinite set formed. This is the regime in which the
reflective theory strictly improves on axiomatic $\mathsf{FM}$ --- both of the
moves that define $\mathsf{FM}$, foundation and the finite-support cut, are
\emph{derived} from the construction rather than \emph{posited}.

\subsection{Going infinitary, and the return of the cut}
\label{sec:fmmodel:infinitary}

Full $\mathsf{FM}$ also has infinite, finitely-supported sets --- $\Atoms$
itself, or the equivariant powerset $\pow_{\fs}(\Atoms)$ of finitely-supported
sets of atoms --- which the finitary signature cannot build. To reach them one
adds an \emph{infinitary} constructor. The most economical is an $\omega$-ary (or
$\kappa$-ary) embrace that consumes a whole stream rather than a finite prefix,
\[
  \inferrule{\Sigma;\Gamma\vdash S \\ (\sigma,\Sigma') = \mathrm{take}(\Sigma,\omega)}
            {\Sigma';\,\sigma,\Gamma\vdash \mathrm{embrace}_\omega(\sigma)\rcup S}
  \quad\textsc{Infinitary-Embrace},
\]
where $\mathrm{take}(\Sigma,\omega)$ returns the entire stream; equivalently one
adjoins a powerset operator and an infinity axiom and climbs the red cumulative
hierarchy $\mathcal{U}_\alpha(\Atoms)$ through all ordinals. Either way two things
change, and --- this is the point the finitary core conceals --- they change
\emph{together}.

\paragraph{An actual infinity, hence one $\omega$ of risk.}
The lazy stream becomes a completed set; $\Atoms$ itself becomes a red set. By
Remark~\ref{rem:risk} the commitment incurred is exactly black infinity, a single
$\omega$-style axiom. No primitive atoms are reintroduced --- the atoms remain
black hereditarily finite sets built from $\bemp$ --- but the potential infinity
of the finitary core is now actual.

\paragraph{Finite support is no longer free; the cut returns.}
Below $\omega$, finite support was automatic. The infinitary hierarchy is not so
forgiving: it contains sets that have \emph{no} finite support. The canonical
witness is the graph of the stream enumeration,
\[
  E \;=\; \{\,(n,a_n) : n\in\omega\,\},
  \qquad a_n = \text{the $n$-th atom of } S[X],
\]
which is fixed by a permutation $\pi$ only when $\pi$ preserves the enumeration,
i.e. only when $\pi$ is the identity on the infinitely many $a_n$; so no finite
set of atoms supports $E$. This is no accident. The formal theory already tells
us that $S[X]$ ``may be thought of as an ordering on $\Atoms$'', and a total order
on an infinite set of atoms is the textbook example of non-finitely-supported
data. The finitary rules never expose that order --- $\mathrm{take}$ reads only a
finite prefix --- which is precisely why support comes for free below $\omega$.
An infinitary embrace that swallows the whole stream imports the order, and with
it the non-finitely-supported sets.

Consequently, in the infinitary regime one must do exactly what axiomatic
$\mathsf{FM}$ does: cut down to the hereditarily finitely-supported part,
\[
  \mathcal{U}_{\mathsf{FM}} \;=\; \{\, x : x \text{ is hereditarily finitely supported}\,\},
\]
discarding $E$ and its kind. On the cut, the some/any theorem and the
equivariance principle hold verbatim, and the result satisfies the Gabbay--Pitts
axioms. The instructive point is that the cut is now the \emph{only} thing
borrowed from axiomatic $\mathsf{FM}$: the atoms still come from the dual copy,
the $\Sym(\Atoms)$-action still comes from $\mathrm{swap}$, and foundation still
comes from the construction (Section~\ref{sec:nwf}); only the finite-support
axiom must be re-imposed, and only because we elected to admit actual infinities.

This locates precisely where reflection saves risk and where it cannot. Atoms:
always free, from the dual copy. Foundation: always free, from local
well-foundedness (Section~\ref{sec:nwf}). Finite support: free in the finitary
core, an imposed cut once infinities are admitted. The finitary/infinitary
boundary is therefore not a technicality but the exact ledger line between what
the construction \emph{derives} and what it must still \emph{assume} --- and a
deliberate design choice for any theory built this way. We would keep the
finitary signature as the canonical core, where the accounting is cleanest, and
treat \textsc{Infinitary-Embrace} as an optional, clearly-priced extension.
).
%% rset.colors.tex
%% Red/black colour layer for the two intertwined copies, plus the FM symbols
%% used by rset.mainthm.tex and rset.applications.tex.
%%
%% EVERYTHING routes through \rd (red copy = player) and \bk (black copy =
%% opponent): redefine just those two macros to change the whole scheme (e.g.
%% to make the "black" copy dark blue for legibility).  Every definition is
%% \providecommand-guarded, so it is safe to \input this file from several
%% places and harmless if rset.local already defines these names.
%%
%% Requires \usepackage{xcolor} (already loaded in rset.tex).

% --- the two copies -------------------------------------------------------
\providecommand{\rd}[1]{\textcolor{red}{#1}}     % red copy  (player)
\providecommand{\bk}[1]{\textcolor{black}{#1}}   % black copy (opponent)

% --- theory names and atom supplies --------------------------------------
\providecommand{\rSetX}{\rd{\mathrm{Set}[X]}}    % red   Set[X]
\providecommand{\bSetX}{\bk{\mathrm{Set}[X]}}    % black Set[X]
\providecommand{\rX}{\rd{X}}                      % red   atom supply
\providecommand{\bX}{\bk{X}}                      % black atom supply

% --- empty set ------------------------------------------------------------
\providecommand{\remp}{\rd{\emptyset}}
\providecommand{\bemp}{\bk{\emptyset}}

% --- membership (relations) ----------------------------------------------
\providecommand{\rin}{\mathrel{\rd{\in}}}
\providecommand{\bin}{\mathrel{\bk{\in}}}
\providecommand{\rnin}{\mathrel{\rd{\notin}}}
\providecommand{\bnin}{\mathrel{\bk{\notin}}}

% --- binary operations ----------------------------------------------------
\providecommand{\rcup}{\mathbin{\rd{\cup}}}
\providecommand{\bcup}{\mathbin{\bk{\cup}}}
\providecommand{\rcap}{\mathbin{\rd{\cap}}}
\providecommand{\bcap}{\mathbin{\bk{\cap}}}
\providecommand{\rsm}{\mathbin{\rd{\setminus}}}
\providecommand{\bsm}{\mathbin{\bk{\setminus}}}

% --- comprehension braces of each theory ---------------------------------
\providecommand{\rbr}[1]{\rd{\{}\,#1\,\rd{\}}}
\providecommand{\bbr}[1]{\bk{\{}\,#1\,\bk{\}}}

% --- colour-blind membership (sees through the colour barrier) ------------
\providecommand{\cbin}{\mathrel{\in_{\!\star}}}

% --- FM symbols (shared with rset.stdfmset.tex) --------------------------
\providecommand{\Atoms}{\mathbb{A}}
\providecommand{\Sym}{\mathrm{Sym}}
\providecommand{\pow}{\mathcal{P}}
\providecommand{\supp}{\mathrm{supp}}
\providecommand{\fs}{\mathrm{fs}}
\providecommand{\new}{\mathbin{\reflectbox{\ensuremath{\mathsf{N}}}}}
\providecommand{\atomsof}{\mathrm{atoms}}
\providecommand{\Red}{\rd{\mathsf{R}}}           % the red HF universe
   % red/black colour layer + FM symbols (providecommand-guarded)

\section{Applications: nominal techniques founded on reflection}
\label{sec:applications}

The original motivation for $\mathsf{FM}$ in computer science was Pitts and
Gabbay's \emph{nominal} treatment of names and binding
\cite{DBLP:journals/fac/GabbayP02}. That treatment is founded squarely on
$\mathsf{FM}$: names are the atoms, $\alpha$-equivalence is an orbit relation
under the permutation action, and the freshness machinery is exactly support and
the $\new$-quantifier. Since Section~\ref{sec:fmmodel} realises $\mathsf{FM}$
inside the reflective universe, the entire nominal edifice descends with it. The
names that nominal techniques take as primitive become, here, black sets built
from the empty set; the reflective construction \emph{founds} the nominal theory
rather than presupposing its atoms.

\subsection{Nominal sets, internally}
\label{sec:applications:nom}

A \emph{nominal set} is a set carrying a $\Sym(\Atoms)$-action in which every
element is finitely supported; nominal sets and equivariant maps form the
category $\mathbf{Nom}$, equivalent to the Schanuel topos. The model
$\mathcal{M}$ of Theorem~\ref{thm:rep} is, on the nose, a nominal set: the action
is that of Section~\ref{sec:fmmodel:action}, and finite support is
Corollary~\ref{cor:hfs}. Every construction nominal practitioners rely on ---
the equivariance principle (anything definable from equivariant data is
equivariant), least supports, and the some/any law for $\new$ --- is therefore
available inside the reflective universe, with the colour discipline of
Section~\ref{sec:fmmodel} keeping the bookkeeping honest: names live in the
{\bk black} copy, the sets that bind them in the {\rd red} copy.

\subsection{Abstract syntax with binding}
\label{sec:applications:syntax}

The headline application of \cite{DBLP:journals/fac/GabbayP02} is a uniform
account of inductively-defined syntax \emph{modulo} $\alpha$-equivalence. Its key
device is \emph{name abstraction}: for a nominal set $X$,
\[
  [\Atoms]X \;=\; (\Atoms\times X)\big/\!\sim,
  \qquad
  \langle a\rangle x \sim \langle b\rangle y
  \;\;\Longleftrightarrow\;\;
  \new c.\ (a\,c)\cdot x = (b\,c)\cdot y,
\]
where $\langle a\rangle x$ is the class of $(a,x)$ and $(a\,c)$ is a
transposition. Both the action and $\new$ are present in $\mathcal{M}$, so
$[\Atoms]X$ is definable there, and it has the property that makes it a model of
binding: $\supp(\langle a\rangle x) = \supp(x)\setminus\{a\}$, i.e. abstraction
\emph{discharges} the bound name from the support, precisely as a binder removes
a variable from the free-variable set.

Inductive syntax is then obtained as usual. Object-level $\lambda$-terms, for
instance, are the least solution of
\[
  \Lambda \;\cong\; \Atoms \;+\; \Lambda\times\Lambda \;+\; [\Atoms]\Lambda
  \qquad(\text{variable, application, abstraction}),
\]
with $\alpha$-equivalence built into the third summand. Structural recursion and
induction come with the \emph{freshness condition for binders}, justified by the
some/any law: a definition or proof that treats a bound name as fresh for its
parameters is sound because, $\Atoms$ being infinite, ``some fresh name'' and
``every fresh name'' agree (Corollary~\ref{cor:hfs}). This is exactly the
reasoning principle nominal users invoke when they ``choose a fresh variable''.

The risk accounting carries through and sharpens the point. Each individual
term --- each element of $\Lambda$ --- is a hereditarily finite nominal set, so
it already lives in the $\mathrm{HF}$ model $\mathcal{M}$ of
Theorem~\ref{thm:rep}, founded on the empty set alone: a binding tree of any
fixed shape costs no infinite risk. The completed datatype $\Lambda$, an
infinite equivariant set, is one of the finitely-supported infinities of
Section~\ref{sec:fmmodel:infinitary}, and so costs exactly one $\omega$ --- black
infinity --- and still not a single primitive name. Where Pitts and Gabbay begin
with an infinite set of atoms given by fiat, here the names are manufactured, the
$\alpha$-structure is an orbit of the $\mathrm{swap}$ action, and the whole
apparatus is refinanced down to one empty set and one successor.

\subsection{Nominal algebraic theories}
\label{sec:applications:alg}

Beyond syntax, equational reasoning in the presence of names --- nominal
equational logic and the nominal Lawvere theories of Clouston
\cite{DBLP:journals/jcss/Clouston14} --- takes its semantics in $\mathbf{Nom}$.
As $\mathcal{M}$ furnishes $\mathbf{Nom}$-structure internally, any such theory
(name-restriction, scope extrusion, binding signatures and their equational
laws) can be interpreted in the reflective universe. The reflective setting adds
one structural feature these accounts do not have: by
Corollary~\ref{cor:two} there are \emph{two} interlocking name-universes, the red
and the black, each supplying the other's names. Reading red as player and black
as opponent recovers the game-semantic picture of the \emph{Intuitions}: a name
on one side is a strategy on the other, and a nominal theory and its dual coexist
in a single reflective universe.

\subsection{Summary}
\label{sec:applications:summary}

The chain is short and worth stating plainly. Nominal techniques rest on
$\mathsf{FM}$; $\mathsf{FM}$ embeds in the reflective universe
(Section~\ref{sec:fmmodel}); the reflective universe is built from the empty set.
Composing, the Pitts--Gabbay theory of names and binding is founded on two
recursively intertwined copies of ordinary set theory, with its primitive atoms
replaced by constructed black sets and its infinite risk replaced by the finite
risk of $\varnothing$ --- exactly the reconciliation announced in the
introduction, now delivered for the application that made $\mathsf{FM}$ matter to
computer science in the first place.
 (add the hook to rset.tex after
%% %% rset.mainthm.tex
%% A model of FM set theory inside the reflective theory of Sections 2--3,
%% with the two copies colour-coded (red = player, black = opponent).
%% Drop in via %% rset.mainthm.tex
%% A model of FM set theory inside the reflective theory of Sections 2--3,
%% with the two copies colour-coded (red = player, black = opponent).
%% Drop in via %% rset.mainthm.tex
%% A model of FM set theory inside the reflective theory of Sections 2--3,
%% with the two copies colour-coded (red = player, black = opponent).
%% Drop in via \input{rset.mainthm} (the hook already exists in rset.tex).
\input{rset.colors}   % red/black colour layer + FM symbols (providecommand-guarded)

\section{A model of $\mathsf{FM}$ in the reflective universe}
\label{sec:fmmodel}

We can now collect the dividend of the construction \emph{Tying the second
recursive knot} above. The reflective theory was assembled from the empty set
alone, yet --- we claim --- it already contains a model of $\mathsf{FM}$, and in
fact two of them, interlocked. The atoms such a model needs are not postulated;
they are \emph{supplied by the dual copy}. Throughout we fix the two tied
theories introduced above, the red copy $\rSetX$ and the black copy $\bSetX$,
with
\[
  \rX \;=\; \bSetX, \qquad \bX \;=\; \rSetX,
\]
so that the admissible atoms of the red theory are exactly the black sets, and
\emph{vice versa}. We write the red copy in {\rd red} and the black copy in
{\bk black} throughout, so that a red comprehension $\rbr{\cdots}$ and a black
comprehension $\bbr{\cdots}$, or the red and black membership relations
$\rin$ and $\bin$, are told apart at a glance.

\subsection{Atoms out of the empty set}
\label{sec:fmmodel:atoms}

Inside the black copy, build the von Neumann numerals from the black empty set,
\[
  \bk{0} \;=\; \bemp,
  \qquad
  \bk{n+1} \;=\; \bk{n} \bcup \bbr{\bk{n}},
\]
each a black, hereditarily finite set obtained, ultimately, from $\bemp$. By
black extensionality (the equations of the formal theory above) these are
pairwise distinct: $\bk{n} \neq \bk{m}$ whenever $n\neq m$. Put
\[
  \Atoms \;:=\; \{\, \bk{n} \;:\; n\in\omega \,\}\ \subseteq\ \bSetX .
\]

Two observations make $\Atoms$ behave \emph{as a set of atoms} for the red
theory.
\begin{itemize}
\item \textbf{Admissibility.} By the second recursive knot $\rX=\bSetX$, every
  black set is an admissible red atom; in particular each $\bk{n}$ is one.
\item \textbf{Opacity.} The red membership relation $\rin$ cannot see into a
  black set (the \emph{Intuitions} above). Hence, \emph{as far as the red theory
  can tell}, an element $\bk{n}$ of $\Atoms$ has no internal structure
  whatsoever: the only facts the red theory may assert about $\bk{n}$ and
  $\bk{m}$ are $\bk{n} = \bk{m}$ or $\bk{n} \neq \bk{m}$, and those are settled,
  out of sight, by black equality.
\end{itemize}
Thus $\Atoms$ presents to the red theory precisely as $\mathsf{FM}$'s set of
atoms does: a countable family of mutually distinct, structureless objects that
may sit inside sets but admit no membership query of their own.

\begin{remark}[Where the risk went]
\label{rem:risk}
This is the sense, promised in the introduction, in which atoms cost no infinite
risk. $\mathsf{ZFA}$ and $\mathsf{FM}$ \emph{postulate} an infinite set of
primitive atoms; the postulate is the risk. Here nothing is postulated beyond
$\bemp$. The atoms are black hereditarily finite sets, and the stream $S[X]$ of
the formal theory hands them out lazily, $\mathrm{take}(\Sigma,n)$ at a time, so
that any red set ever built mentions only finitely many of them. The supply is
\emph{potentially} infinite, never an actual completed infinity: the commitment
is the black empty set together with the same successor step that underwrites
$\omega$, and not one primitive atom more.
\end{remark}

\subsection{The permutation action}
\label{sec:fmmodel:action}

The stream operation $\mathrm{swap}(S[X],\rho)$ of the formal theory applies a
finite permutation $\rho=\{\,\bk{j_1}/\bk{i_1}:\dots:\bk{j_k}/\bk{i_k}\,\}$ to
the ordering of atoms; these finite permutations generate the full group
$\Sym(\Atoms)$. A permutation $\pi\in\Sym(\Atoms)$ acts on red sets by
red-membership recursion exactly as in $\mathsf{FM}$,
\[
  \pi\cdot a = \pi(a)\ (a\in\Atoms),
  \qquad
  \pi\cdot S = \rd{\{}\, \pi\cdot e \;:\; e \rin S \,\rd{\}}.
\]
Because $\pi$ moves only atoms and never the black interior of an atom, and
because the red operations $\remp,\ \rcup,\ \rcap,\ \rsm$ and
$\mathrm{embrace}_i$ are visibly equivariant (a renaming commutes with each), the
action is by automorphisms of the red theory and respects the equations of the
formal theory. We have, verbatim, the $\Sym(\Atoms)$-action that $\mathsf{FM}$
requires.

\subsection{Finite support is automatic}
\label{sec:fmmodel:support}

Write $\Red$ for the class of well-formed red sets, taken modulo red equality,
all of whose atoms lie in $\Atoms$. The finitary rules of the formal theory ---
$\mathrm{embrace}_i$ has finite arity, $\mathrm{Atomicity}$ adjoins one atom, and
$\rcup,\rcap,\rsm$ and $\mathrm{Comprehension}$ over finite sets preserve
finiteness --- yield only \emph{hereditarily finite} red sets. For $S\in\Red$ let
$\atomsof(S)\subseteq\Atoms$ be the finite set of atoms occurring hereditarily in
$S$.

\begin{lemma}[Support]
\label{lem:support}
For every $S\in\Red$ the finite set $\atomsof(S)$ supports $S$: if
$\pi\in\Sym(\Atoms)$ fixes $\atomsof(S)$ pointwise then $\pi\cdot S = S$.
Moreover $\atomsof(S)$ is the \emph{least} support, so $\supp(S)=\atomsof(S)$.
\end{lemma}

\begin{proof}[Proof sketch]
That $\atomsof(S)$ supports $S$ is an induction on the rank of $S$: an atom
$a\in\atomsof(S)$ is fixed by hypothesis, and for a set $S=\rbr{e_1,\dots,e_m}$ a
permutation fixing $\atomsof(S)\supseteq\bigcup_i \atomsof(e_i)$ fixes each $e_i$
by induction, hence fixes $S$. Leastness uses that $\Atoms$ is infinite: if
$a\in\atomsof(S)$ then, choosing a fresh $b\notin\atomsof(S)$, the transposition
$(a\,b)$ fixes $\atomsof(S)\setminus\{a\}$ yet moves $S$, so $a$ lies in every
support. Hence no proper subset of $\atomsof(S)$ supports $S$.
\end{proof}

\begin{corollary}
\label{cor:hfs}
Every $S\in\Red$ is hereditarily finitely supported, with
$\supp(S)=\atomsof(S)$. Freshness is occurrence-failure,
$a\mathrel{\#}S \iff a\notin\atomsof(S)$, and since $\Atoms$ is infinite the
fresh quantifier $\new$ and its some/any law (Section~\ref{sec:fm}) hold in
$\Red$.
\end{corollary}

The point worth stressing is that in $\mathsf{FM}$ finite support is an
\emph{axiom}, imposed by cutting the cumulative hierarchy down to
$\mathcal{U}_{\mathsf{FM}}$. Here it is a \emph{theorem}: the finitary red rules
never build anything that is not finitely supported, so the cut is vacuous and
nothing need be assumed.

\begin{remark}[Two points of precision]
\label{rem:group}
Lemma~\ref{lem:support} rests on two facts worth making explicit. First, the
group in play is the \emph{finitary} symmetric group of $\Atoms$ --- the
permutations generated by the transpositions that $\mathrm{swap}$ realises --- for
which supports are closed under intersection and least supports therefore exist;
on an arbitrary subgroup they need not. Second, the action permutes atoms while
treating each as structureless: it is an automorphism of the \emph{red} structure
$(\Red,\rin)$ alone, and is not asked to respect the black membership $\bin$
internal to an atom. The atoms are accordingly black-distinct but red-opaque, and
the entire support theory lives at the red level.
\end{remark}

\subsection{The representation theorem}
\label{sec:fmmodel:thm}

Let $\mathcal{M}$ be the structure with universe $\Atoms\cup\Red$, with the atoms
$\Atoms$ as urelements, with $\in^{\mathcal M}$ the red membership relation and
$=^{\mathcal M}$ red equality.

\begin{theorem}[Representation]
\label{thm:rep}
$\mathcal{M}$ is a model of hereditarily-finite Fraenkel--Mostowski set theory:
the $\mathsf{FM}$ analogue of $\mathrm{HF}(\Atoms)$. Concretely, interpreting the
atoms of $\mathsf{FM}$ as $\Atoms$, its sets as $\Red$, and its membership as red
$\rin$, the structure $\mathcal{M}$ validates extensionality, the empty set,
pairing, (finite) union, (finite) separation and foundation in their
$\mathsf{ZFA}$ form, carries the $\Sym(\Atoms)$-action of
Section~\ref{sec:fmmodel:action}, and satisfies the finite-support axiom. The
fresh quantifier and the some/any theorem hold throughout.
\end{theorem}

\begin{proof}[Proof sketch]
The atoms are membersless: red $\rin$ cannot enter a black set, so no
$a\in\Atoms$ has red elements, and atoms are $\in$-minimal and disjoint from
$\Red$ in sort, as $\mathsf{ZFA}$ demands. Foundation holds because each red set
is built at a finite rank over the atoms, so $\in^{\mathcal M}$ is well-founded.
The remaining axioms come straight from the rules of the formal theory:
$\mathrm{Foundation}$ gives $\remp$; $\mathrm{embrace}$ with $\mathrm{Union}$
gives pairing $\rbr{e_1}\rcup\rbr{e_2}$ and, iterated, every finite set;
$\mathrm{Union}$ gives binary and hence finite unions; $\mathrm{Subtraction}$ and
$\mathrm{Comprehension}$ over finite domains give finite separation; and
extensionality holds for the model provided provable red equality coincides with
sameness of members --- a property the equations of the formal theory are intended
to force, and which Section~\ref{sec:rigor} records as one of the obligations a
fully rigorous development must discharge. The $\Sym(\Atoms)$-action and the
finite-support axiom are Section~\ref{sec:fmmodel:action} and
Corollary~\ref{cor:hfs}; the some/any law follows because $\Atoms$ is infinite
and every $S\in\Red$ has finite support. Powerset and infinity are absent --- as
expected of an $\mathrm{HF}$ universe --- and are the subject of
Section~\ref{sec:fmmodel:infinitary}.
\end{proof}

\begin{corollary}[Two interlocked models]
\label{cor:two}
By the red/black symmetry of the construction, the same argument run with the
colours exchanged produces a black model $\bk{\mathcal M}$ of
hereditarily-finite $\mathsf{FM}$ whose atoms are the red numerals
$\{\,\rd{n} : n\in\omega\,\}\subseteq\rSetX$. The two models interlock: the atoms
of each are sets of the other. In the game-semantic reading of the
\emph{Intuitions} above, player's atoms are opponent's strategies and conversely
--- a single reflective universe carrying both $\mathsf{FM}$ models at once, each
underwriting the other's names.
\end{corollary}

\begin{remark}[What the $\mathrm{HF}$ restriction costs]
\label{rem:hftrivial}
It should be said plainly what is and is not delivered. The phenomena that give
$\mathsf{FM}$ its character --- finite support as a genuine constraint on
\emph{infinite} objects, equivariance with non-trivial orbits, and the some/any
law quantifying over infinitely many fresh names against a fixed infinite
support --- all require infinite, finitely-supported sets, of which $\mathcal{M}$
has none. On the hereditarily finite fragment these properties hold, but
trivially: every set is finite, so finite support is automatic and the some/any
law degenerates. $\mathcal{M}$ is thus best described as a model of the
\emph{nominal hereditarily finite sets}: faithful to the $\mathsf{ZFA}$ and
$\mathsf{FM}$ axioms, but exercising the distinctively nominal content only once
the infinitary extension of Section~\ref{sec:fmmodel:infinitary} is in force. We
retain the phrase ``model of $\mathsf{FM}$'' only under this qualification.
\end{remark}

\subsection{Finitary by design: the hereditarily finite core}
\label{sec:fmmodel:finitary}

It is worth being explicit about why Theorem~\ref{thm:rep} lands on the
hereditarily finite fragment, because the reason is also the source of the
construction's advantage. The signature of the formal theory offers, as set
constructors, only $\remp$, the finite-arity embraces $\mathrm{embrace}_i$
($|\mathrm{embrace}_i| = i < \omega$), the binary $\rcup,\rcap,\rsm$, and
$\mathrm{Comprehension}$ over finite domains $S_1\times\dots\times S_n$. Each of
these maps finite sets to finite sets, and the only base case is $\remp$;
so a trivial induction shows that \emph{every} derivable red set is hereditarily
finite. The model is $\mathrm{HF}(\Atoms)$, the nominal hereditarily finite sets
over the atoms.

Two of $\mathsf{FM}$'s defining commitments are, in this regime, theorems rather
than axioms.
\begin{itemize}
\item \textbf{Foundation} holds because every red set sits at a finite rank over
  the atoms, so $\in^{\mathcal{M}}$ is well-founded (Theorem~\ref{thm:rep}); we
  never assumed regularity.
\item \textbf{Finite support} holds because a hereditarily finite set mentions
  only finitely many atoms, each of which is its own least support
  (Lemma~\ref{lem:support}); we never assumed the finite-support axiom.
\end{itemize}
Correspondingly there is no infinity axiom and no infinite risk: the stream
$S[X]$ is read only through $\mathrm{take}(\Sigma,n)$ for finite $n$, so at no
point is an actually-infinite set formed. This is the regime in which the
reflective theory strictly improves on axiomatic $\mathsf{FM}$ --- both of the
moves that define $\mathsf{FM}$, foundation and the finite-support cut, are
\emph{derived} from the construction rather than \emph{posited}.

\subsection{Going infinitary, and the return of the cut}
\label{sec:fmmodel:infinitary}

Full $\mathsf{FM}$ also has infinite, finitely-supported sets --- $\Atoms$
itself, or the equivariant powerset $\pow_{\fs}(\Atoms)$ of finitely-supported
sets of atoms --- which the finitary signature cannot build. To reach them one
adds an \emph{infinitary} constructor. The most economical is an $\omega$-ary (or
$\kappa$-ary) embrace that consumes a whole stream rather than a finite prefix,
\[
  \inferrule{\Sigma;\Gamma\vdash S \\ (\sigma,\Sigma') = \mathrm{take}(\Sigma,\omega)}
            {\Sigma';\,\sigma,\Gamma\vdash \mathrm{embrace}_\omega(\sigma)\rcup S}
  \quad\textsc{Infinitary-Embrace},
\]
where $\mathrm{take}(\Sigma,\omega)$ returns the entire stream; equivalently one
adjoins a powerset operator and an infinity axiom and climbs the red cumulative
hierarchy $\mathcal{U}_\alpha(\Atoms)$ through all ordinals. Either way two things
change, and --- this is the point the finitary core conceals --- they change
\emph{together}.

\paragraph{An actual infinity, hence one $\omega$ of risk.}
The lazy stream becomes a completed set; $\Atoms$ itself becomes a red set. By
Remark~\ref{rem:risk} the commitment incurred is exactly black infinity, a single
$\omega$-style axiom. No primitive atoms are reintroduced --- the atoms remain
black hereditarily finite sets built from $\bemp$ --- but the potential infinity
of the finitary core is now actual.

\paragraph{Finite support is no longer free; the cut returns.}
Below $\omega$, finite support was automatic. The infinitary hierarchy is not so
forgiving: it contains sets that have \emph{no} finite support. The canonical
witness is the graph of the stream enumeration,
\[
  E \;=\; \{\,(n,a_n) : n\in\omega\,\},
  \qquad a_n = \text{the $n$-th atom of } S[X],
\]
which is fixed by a permutation $\pi$ only when $\pi$ preserves the enumeration,
i.e. only when $\pi$ is the identity on the infinitely many $a_n$; so no finite
set of atoms supports $E$. This is no accident. The formal theory already tells
us that $S[X]$ ``may be thought of as an ordering on $\Atoms$'', and a total order
on an infinite set of atoms is the textbook example of non-finitely-supported
data. The finitary rules never expose that order --- $\mathrm{take}$ reads only a
finite prefix --- which is precisely why support comes for free below $\omega$.
An infinitary embrace that swallows the whole stream imports the order, and with
it the non-finitely-supported sets.

Consequently, in the infinitary regime one must do exactly what axiomatic
$\mathsf{FM}$ does: cut down to the hereditarily finitely-supported part,
\[
  \mathcal{U}_{\mathsf{FM}} \;=\; \{\, x : x \text{ is hereditarily finitely supported}\,\},
\]
discarding $E$ and its kind. On the cut, the some/any theorem and the
equivariance principle hold verbatim, and the result satisfies the Gabbay--Pitts
axioms. The instructive point is that the cut is now the \emph{only} thing
borrowed from axiomatic $\mathsf{FM}$: the atoms still come from the dual copy,
the $\Sym(\Atoms)$-action still comes from $\mathrm{swap}$, and foundation still
comes from the construction (Section~\ref{sec:nwf}); only the finite-support
axiom must be re-imposed, and only because we elected to admit actual infinities.

This locates precisely where reflection saves risk and where it cannot. Atoms:
always free, from the dual copy. Foundation: always free, from local
well-foundedness (Section~\ref{sec:nwf}). Finite support: free in the finitary
core, an imposed cut once infinities are admitted. The finitary/infinitary
boundary is therefore not a technicality but the exact ledger line between what
the construction \emph{derives} and what it must still \emph{assume} --- and a
deliberate design choice for any theory built this way. We would keep the
finitary signature as the canonical core, where the accounting is cleanest, and
treat \textsc{Infinitary-Embrace} as an optional, clearly-priced extension.
 (the hook already exists in rset.tex).
%% rset.colors.tex
%% Red/black colour layer for the two intertwined copies, plus the FM symbols
%% used by rset.mainthm.tex and rset.applications.tex.
%%
%% EVERYTHING routes through \rd (red copy = player) and \bk (black copy =
%% opponent): redefine just those two macros to change the whole scheme (e.g.
%% to make the "black" copy dark blue for legibility).  Every definition is
%% \providecommand-guarded, so it is safe to \input this file from several
%% places and harmless if rset.local already defines these names.
%%
%% Requires \usepackage{xcolor} (already loaded in rset.tex).

% --- the two copies -------------------------------------------------------
\providecommand{\rd}[1]{\textcolor{red}{#1}}     % red copy  (player)
\providecommand{\bk}[1]{\textcolor{black}{#1}}   % black copy (opponent)

% --- theory names and atom supplies --------------------------------------
\providecommand{\rSetX}{\rd{\mathrm{Set}[X]}}    % red   Set[X]
\providecommand{\bSetX}{\bk{\mathrm{Set}[X]}}    % black Set[X]
\providecommand{\rX}{\rd{X}}                      % red   atom supply
\providecommand{\bX}{\bk{X}}                      % black atom supply

% --- empty set ------------------------------------------------------------
\providecommand{\remp}{\rd{\emptyset}}
\providecommand{\bemp}{\bk{\emptyset}}

% --- membership (relations) ----------------------------------------------
\providecommand{\rin}{\mathrel{\rd{\in}}}
\providecommand{\bin}{\mathrel{\bk{\in}}}
\providecommand{\rnin}{\mathrel{\rd{\notin}}}
\providecommand{\bnin}{\mathrel{\bk{\notin}}}

% --- binary operations ----------------------------------------------------
\providecommand{\rcup}{\mathbin{\rd{\cup}}}
\providecommand{\bcup}{\mathbin{\bk{\cup}}}
\providecommand{\rcap}{\mathbin{\rd{\cap}}}
\providecommand{\bcap}{\mathbin{\bk{\cap}}}
\providecommand{\rsm}{\mathbin{\rd{\setminus}}}
\providecommand{\bsm}{\mathbin{\bk{\setminus}}}

% --- comprehension braces of each theory ---------------------------------
\providecommand{\rbr}[1]{\rd{\{}\,#1\,\rd{\}}}
\providecommand{\bbr}[1]{\bk{\{}\,#1\,\bk{\}}}

% --- colour-blind membership (sees through the colour barrier) ------------
\providecommand{\cbin}{\mathrel{\in_{\!\star}}}

% --- FM symbols (shared with rset.stdfmset.tex) --------------------------
\providecommand{\Atoms}{\mathbb{A}}
\providecommand{\Sym}{\mathrm{Sym}}
\providecommand{\pow}{\mathcal{P}}
\providecommand{\supp}{\mathrm{supp}}
\providecommand{\fs}{\mathrm{fs}}
\providecommand{\new}{\mathbin{\reflectbox{\ensuremath{\mathsf{N}}}}}
\providecommand{\atomsof}{\mathrm{atoms}}
\providecommand{\Red}{\rd{\mathsf{R}}}           % the red HF universe
   % red/black colour layer + FM symbols (providecommand-guarded)

\section{A model of $\mathsf{FM}$ in the reflective universe}
\label{sec:fmmodel}

We can now collect the dividend of the construction \emph{Tying the second
recursive knot} above. The reflective theory was assembled from the empty set
alone, yet --- we claim --- it already contains a model of $\mathsf{FM}$, and in
fact two of them, interlocked. The atoms such a model needs are not postulated;
they are \emph{supplied by the dual copy}. Throughout we fix the two tied
theories introduced above, the red copy $\rSetX$ and the black copy $\bSetX$,
with
\[
  \rX \;=\; \bSetX, \qquad \bX \;=\; \rSetX,
\]
so that the admissible atoms of the red theory are exactly the black sets, and
\emph{vice versa}. We write the red copy in {\rd red} and the black copy in
{\bk black} throughout, so that a red comprehension $\rbr{\cdots}$ and a black
comprehension $\bbr{\cdots}$, or the red and black membership relations
$\rin$ and $\bin$, are told apart at a glance.

\subsection{Atoms out of the empty set}
\label{sec:fmmodel:atoms}

Inside the black copy, build the von Neumann numerals from the black empty set,
\[
  \bk{0} \;=\; \bemp,
  \qquad
  \bk{n+1} \;=\; \bk{n} \bcup \bbr{\bk{n}},
\]
each a black, hereditarily finite set obtained, ultimately, from $\bemp$. By
black extensionality (the equations of the formal theory above) these are
pairwise distinct: $\bk{n} \neq \bk{m}$ whenever $n\neq m$. Put
\[
  \Atoms \;:=\; \{\, \bk{n} \;:\; n\in\omega \,\}\ \subseteq\ \bSetX .
\]

Two observations make $\Atoms$ behave \emph{as a set of atoms} for the red
theory.
\begin{itemize}
\item \textbf{Admissibility.} By the second recursive knot $\rX=\bSetX$, every
  black set is an admissible red atom; in particular each $\bk{n}$ is one.
\item \textbf{Opacity.} The red membership relation $\rin$ cannot see into a
  black set (the \emph{Intuitions} above). Hence, \emph{as far as the red theory
  can tell}, an element $\bk{n}$ of $\Atoms$ has no internal structure
  whatsoever: the only facts the red theory may assert about $\bk{n}$ and
  $\bk{m}$ are $\bk{n} = \bk{m}$ or $\bk{n} \neq \bk{m}$, and those are settled,
  out of sight, by black equality.
\end{itemize}
Thus $\Atoms$ presents to the red theory precisely as $\mathsf{FM}$'s set of
atoms does: a countable family of mutually distinct, structureless objects that
may sit inside sets but admit no membership query of their own.

\begin{remark}[Where the risk went]
\label{rem:risk}
This is the sense, promised in the introduction, in which atoms cost no infinite
risk. $\mathsf{ZFA}$ and $\mathsf{FM}$ \emph{postulate} an infinite set of
primitive atoms; the postulate is the risk. Here nothing is postulated beyond
$\bemp$. The atoms are black hereditarily finite sets, and the stream $S[X]$ of
the formal theory hands them out lazily, $\mathrm{take}(\Sigma,n)$ at a time, so
that any red set ever built mentions only finitely many of them. The supply is
\emph{potentially} infinite, never an actual completed infinity: the commitment
is the black empty set together with the same successor step that underwrites
$\omega$, and not one primitive atom more.
\end{remark}

\subsection{The permutation action}
\label{sec:fmmodel:action}

The stream operation $\mathrm{swap}(S[X],\rho)$ of the formal theory applies a
finite permutation $\rho=\{\,\bk{j_1}/\bk{i_1}:\dots:\bk{j_k}/\bk{i_k}\,\}$ to
the ordering of atoms; these finite permutations generate the full group
$\Sym(\Atoms)$. A permutation $\pi\in\Sym(\Atoms)$ acts on red sets by
red-membership recursion exactly as in $\mathsf{FM}$,
\[
  \pi\cdot a = \pi(a)\ (a\in\Atoms),
  \qquad
  \pi\cdot S = \rd{\{}\, \pi\cdot e \;:\; e \rin S \,\rd{\}}.
\]
Because $\pi$ moves only atoms and never the black interior of an atom, and
because the red operations $\remp,\ \rcup,\ \rcap,\ \rsm$ and
$\mathrm{embrace}_i$ are visibly equivariant (a renaming commutes with each), the
action is by automorphisms of the red theory and respects the equations of the
formal theory. We have, verbatim, the $\Sym(\Atoms)$-action that $\mathsf{FM}$
requires.

\subsection{Finite support is automatic}
\label{sec:fmmodel:support}

Write $\Red$ for the class of well-formed red sets, taken modulo red equality,
all of whose atoms lie in $\Atoms$. The finitary rules of the formal theory ---
$\mathrm{embrace}_i$ has finite arity, $\mathrm{Atomicity}$ adjoins one atom, and
$\rcup,\rcap,\rsm$ and $\mathrm{Comprehension}$ over finite sets preserve
finiteness --- yield only \emph{hereditarily finite} red sets. For $S\in\Red$ let
$\atomsof(S)\subseteq\Atoms$ be the finite set of atoms occurring hereditarily in
$S$.

\begin{lemma}[Support]
\label{lem:support}
For every $S\in\Red$ the finite set $\atomsof(S)$ supports $S$: if
$\pi\in\Sym(\Atoms)$ fixes $\atomsof(S)$ pointwise then $\pi\cdot S = S$.
Moreover $\atomsof(S)$ is the \emph{least} support, so $\supp(S)=\atomsof(S)$.
\end{lemma}

\begin{proof}[Proof sketch]
That $\atomsof(S)$ supports $S$ is an induction on the rank of $S$: an atom
$a\in\atomsof(S)$ is fixed by hypothesis, and for a set $S=\rbr{e_1,\dots,e_m}$ a
permutation fixing $\atomsof(S)\supseteq\bigcup_i \atomsof(e_i)$ fixes each $e_i$
by induction, hence fixes $S$. Leastness uses that $\Atoms$ is infinite: if
$a\in\atomsof(S)$ then, choosing a fresh $b\notin\atomsof(S)$, the transposition
$(a\,b)$ fixes $\atomsof(S)\setminus\{a\}$ yet moves $S$, so $a$ lies in every
support. Hence no proper subset of $\atomsof(S)$ supports $S$.
\end{proof}

\begin{corollary}
\label{cor:hfs}
Every $S\in\Red$ is hereditarily finitely supported, with
$\supp(S)=\atomsof(S)$. Freshness is occurrence-failure,
$a\mathrel{\#}S \iff a\notin\atomsof(S)$, and since $\Atoms$ is infinite the
fresh quantifier $\new$ and its some/any law (Section~\ref{sec:fm}) hold in
$\Red$.
\end{corollary}

The point worth stressing is that in $\mathsf{FM}$ finite support is an
\emph{axiom}, imposed by cutting the cumulative hierarchy down to
$\mathcal{U}_{\mathsf{FM}}$. Here it is a \emph{theorem}: the finitary red rules
never build anything that is not finitely supported, so the cut is vacuous and
nothing need be assumed.

\begin{remark}[Two points of precision]
\label{rem:group}
Lemma~\ref{lem:support} rests on two facts worth making explicit. First, the
group in play is the \emph{finitary} symmetric group of $\Atoms$ --- the
permutations generated by the transpositions that $\mathrm{swap}$ realises --- for
which supports are closed under intersection and least supports therefore exist;
on an arbitrary subgroup they need not. Second, the action permutes atoms while
treating each as structureless: it is an automorphism of the \emph{red} structure
$(\Red,\rin)$ alone, and is not asked to respect the black membership $\bin$
internal to an atom. The atoms are accordingly black-distinct but red-opaque, and
the entire support theory lives at the red level.
\end{remark}

\subsection{The representation theorem}
\label{sec:fmmodel:thm}

Let $\mathcal{M}$ be the structure with universe $\Atoms\cup\Red$, with the atoms
$\Atoms$ as urelements, with $\in^{\mathcal M}$ the red membership relation and
$=^{\mathcal M}$ red equality.

\begin{theorem}[Representation]
\label{thm:rep}
$\mathcal{M}$ is a model of hereditarily-finite Fraenkel--Mostowski set theory:
the $\mathsf{FM}$ analogue of $\mathrm{HF}(\Atoms)$. Concretely, interpreting the
atoms of $\mathsf{FM}$ as $\Atoms$, its sets as $\Red$, and its membership as red
$\rin$, the structure $\mathcal{M}$ validates extensionality, the empty set,
pairing, (finite) union, (finite) separation and foundation in their
$\mathsf{ZFA}$ form, carries the $\Sym(\Atoms)$-action of
Section~\ref{sec:fmmodel:action}, and satisfies the finite-support axiom. The
fresh quantifier and the some/any theorem hold throughout.
\end{theorem}

\begin{proof}[Proof sketch]
The atoms are membersless: red $\rin$ cannot enter a black set, so no
$a\in\Atoms$ has red elements, and atoms are $\in$-minimal and disjoint from
$\Red$ in sort, as $\mathsf{ZFA}$ demands. Foundation holds because each red set
is built at a finite rank over the atoms, so $\in^{\mathcal M}$ is well-founded.
The remaining axioms come straight from the rules of the formal theory:
$\mathrm{Foundation}$ gives $\remp$; $\mathrm{embrace}$ with $\mathrm{Union}$
gives pairing $\rbr{e_1}\rcup\rbr{e_2}$ and, iterated, every finite set;
$\mathrm{Union}$ gives binary and hence finite unions; $\mathrm{Subtraction}$ and
$\mathrm{Comprehension}$ over finite domains give finite separation; and
extensionality holds for the model provided provable red equality coincides with
sameness of members --- a property the equations of the formal theory are intended
to force, and which Section~\ref{sec:rigor} records as one of the obligations a
fully rigorous development must discharge. The $\Sym(\Atoms)$-action and the
finite-support axiom are Section~\ref{sec:fmmodel:action} and
Corollary~\ref{cor:hfs}; the some/any law follows because $\Atoms$ is infinite
and every $S\in\Red$ has finite support. Powerset and infinity are absent --- as
expected of an $\mathrm{HF}$ universe --- and are the subject of
Section~\ref{sec:fmmodel:infinitary}.
\end{proof}

\begin{corollary}[Two interlocked models]
\label{cor:two}
By the red/black symmetry of the construction, the same argument run with the
colours exchanged produces a black model $\bk{\mathcal M}$ of
hereditarily-finite $\mathsf{FM}$ whose atoms are the red numerals
$\{\,\rd{n} : n\in\omega\,\}\subseteq\rSetX$. The two models interlock: the atoms
of each are sets of the other. In the game-semantic reading of the
\emph{Intuitions} above, player's atoms are opponent's strategies and conversely
--- a single reflective universe carrying both $\mathsf{FM}$ models at once, each
underwriting the other's names.
\end{corollary}

\begin{remark}[What the $\mathrm{HF}$ restriction costs]
\label{rem:hftrivial}
It should be said plainly what is and is not delivered. The phenomena that give
$\mathsf{FM}$ its character --- finite support as a genuine constraint on
\emph{infinite} objects, equivariance with non-trivial orbits, and the some/any
law quantifying over infinitely many fresh names against a fixed infinite
support --- all require infinite, finitely-supported sets, of which $\mathcal{M}$
has none. On the hereditarily finite fragment these properties hold, but
trivially: every set is finite, so finite support is automatic and the some/any
law degenerates. $\mathcal{M}$ is thus best described as a model of the
\emph{nominal hereditarily finite sets}: faithful to the $\mathsf{ZFA}$ and
$\mathsf{FM}$ axioms, but exercising the distinctively nominal content only once
the infinitary extension of Section~\ref{sec:fmmodel:infinitary} is in force. We
retain the phrase ``model of $\mathsf{FM}$'' only under this qualification.
\end{remark}

\subsection{Finitary by design: the hereditarily finite core}
\label{sec:fmmodel:finitary}

It is worth being explicit about why Theorem~\ref{thm:rep} lands on the
hereditarily finite fragment, because the reason is also the source of the
construction's advantage. The signature of the formal theory offers, as set
constructors, only $\remp$, the finite-arity embraces $\mathrm{embrace}_i$
($|\mathrm{embrace}_i| = i < \omega$), the binary $\rcup,\rcap,\rsm$, and
$\mathrm{Comprehension}$ over finite domains $S_1\times\dots\times S_n$. Each of
these maps finite sets to finite sets, and the only base case is $\remp$;
so a trivial induction shows that \emph{every} derivable red set is hereditarily
finite. The model is $\mathrm{HF}(\Atoms)$, the nominal hereditarily finite sets
over the atoms.

Two of $\mathsf{FM}$'s defining commitments are, in this regime, theorems rather
than axioms.
\begin{itemize}
\item \textbf{Foundation} holds because every red set sits at a finite rank over
  the atoms, so $\in^{\mathcal{M}}$ is well-founded (Theorem~\ref{thm:rep}); we
  never assumed regularity.
\item \textbf{Finite support} holds because a hereditarily finite set mentions
  only finitely many atoms, each of which is its own least support
  (Lemma~\ref{lem:support}); we never assumed the finite-support axiom.
\end{itemize}
Correspondingly there is no infinity axiom and no infinite risk: the stream
$S[X]$ is read only through $\mathrm{take}(\Sigma,n)$ for finite $n$, so at no
point is an actually-infinite set formed. This is the regime in which the
reflective theory strictly improves on axiomatic $\mathsf{FM}$ --- both of the
moves that define $\mathsf{FM}$, foundation and the finite-support cut, are
\emph{derived} from the construction rather than \emph{posited}.

\subsection{Going infinitary, and the return of the cut}
\label{sec:fmmodel:infinitary}

Full $\mathsf{FM}$ also has infinite, finitely-supported sets --- $\Atoms$
itself, or the equivariant powerset $\pow_{\fs}(\Atoms)$ of finitely-supported
sets of atoms --- which the finitary signature cannot build. To reach them one
adds an \emph{infinitary} constructor. The most economical is an $\omega$-ary (or
$\kappa$-ary) embrace that consumes a whole stream rather than a finite prefix,
\[
  \inferrule{\Sigma;\Gamma\vdash S \\ (\sigma,\Sigma') = \mathrm{take}(\Sigma,\omega)}
            {\Sigma';\,\sigma,\Gamma\vdash \mathrm{embrace}_\omega(\sigma)\rcup S}
  \quad\textsc{Infinitary-Embrace},
\]
where $\mathrm{take}(\Sigma,\omega)$ returns the entire stream; equivalently one
adjoins a powerset operator and an infinity axiom and climbs the red cumulative
hierarchy $\mathcal{U}_\alpha(\Atoms)$ through all ordinals. Either way two things
change, and --- this is the point the finitary core conceals --- they change
\emph{together}.

\paragraph{An actual infinity, hence one $\omega$ of risk.}
The lazy stream becomes a completed set; $\Atoms$ itself becomes a red set. By
Remark~\ref{rem:risk} the commitment incurred is exactly black infinity, a single
$\omega$-style axiom. No primitive atoms are reintroduced --- the atoms remain
black hereditarily finite sets built from $\bemp$ --- but the potential infinity
of the finitary core is now actual.

\paragraph{Finite support is no longer free; the cut returns.}
Below $\omega$, finite support was automatic. The infinitary hierarchy is not so
forgiving: it contains sets that have \emph{no} finite support. The canonical
witness is the graph of the stream enumeration,
\[
  E \;=\; \{\,(n,a_n) : n\in\omega\,\},
  \qquad a_n = \text{the $n$-th atom of } S[X],
\]
which is fixed by a permutation $\pi$ only when $\pi$ preserves the enumeration,
i.e. only when $\pi$ is the identity on the infinitely many $a_n$; so no finite
set of atoms supports $E$. This is no accident. The formal theory already tells
us that $S[X]$ ``may be thought of as an ordering on $\Atoms$'', and a total order
on an infinite set of atoms is the textbook example of non-finitely-supported
data. The finitary rules never expose that order --- $\mathrm{take}$ reads only a
finite prefix --- which is precisely why support comes for free below $\omega$.
An infinitary embrace that swallows the whole stream imports the order, and with
it the non-finitely-supported sets.

Consequently, in the infinitary regime one must do exactly what axiomatic
$\mathsf{FM}$ does: cut down to the hereditarily finitely-supported part,
\[
  \mathcal{U}_{\mathsf{FM}} \;=\; \{\, x : x \text{ is hereditarily finitely supported}\,\},
\]
discarding $E$ and its kind. On the cut, the some/any theorem and the
equivariance principle hold verbatim, and the result satisfies the Gabbay--Pitts
axioms. The instructive point is that the cut is now the \emph{only} thing
borrowed from axiomatic $\mathsf{FM}$: the atoms still come from the dual copy,
the $\Sym(\Atoms)$-action still comes from $\mathrm{swap}$, and foundation still
comes from the construction (Section~\ref{sec:nwf}); only the finite-support
axiom must be re-imposed, and only because we elected to admit actual infinities.

This locates precisely where reflection saves risk and where it cannot. Atoms:
always free, from the dual copy. Foundation: always free, from local
well-foundedness (Section~\ref{sec:nwf}). Finite support: free in the finitary
core, an imposed cut once infinities are admitted. The finitary/infinitary
boundary is therefore not a technicality but the exact ledger line between what
the construction \emph{derives} and what it must still \emph{assume} --- and a
deliberate design choice for any theory built this way. We would keep the
finitary signature as the canonical core, where the accounting is cleanest, and
treat \textsc{Infinitary-Embrace} as an optional, clearly-priced extension.
 (the hook already exists in rset.tex).
%% rset.colors.tex
%% Red/black colour layer for the two intertwined copies, plus the FM symbols
%% used by rset.mainthm.tex and rset.applications.tex.
%%
%% EVERYTHING routes through \rd (red copy = player) and \bk (black copy =
%% opponent): redefine just those two macros to change the whole scheme (e.g.
%% to make the "black" copy dark blue for legibility).  Every definition is
%% \providecommand-guarded, so it is safe to \input this file from several
%% places and harmless if rset.local already defines these names.
%%
%% Requires \usepackage{xcolor} (already loaded in rset.tex).

% --- the two copies -------------------------------------------------------
\providecommand{\rd}[1]{\textcolor{red}{#1}}     % red copy  (player)
\providecommand{\bk}[1]{\textcolor{black}{#1}}   % black copy (opponent)

% --- theory names and atom supplies --------------------------------------
\providecommand{\rSetX}{\rd{\mathrm{Set}[X]}}    % red   Set[X]
\providecommand{\bSetX}{\bk{\mathrm{Set}[X]}}    % black Set[X]
\providecommand{\rX}{\rd{X}}                      % red   atom supply
\providecommand{\bX}{\bk{X}}                      % black atom supply

% --- empty set ------------------------------------------------------------
\providecommand{\remp}{\rd{\emptyset}}
\providecommand{\bemp}{\bk{\emptyset}}

% --- membership (relations) ----------------------------------------------
\providecommand{\rin}{\mathrel{\rd{\in}}}
\providecommand{\bin}{\mathrel{\bk{\in}}}
\providecommand{\rnin}{\mathrel{\rd{\notin}}}
\providecommand{\bnin}{\mathrel{\bk{\notin}}}

% --- binary operations ----------------------------------------------------
\providecommand{\rcup}{\mathbin{\rd{\cup}}}
\providecommand{\bcup}{\mathbin{\bk{\cup}}}
\providecommand{\rcap}{\mathbin{\rd{\cap}}}
\providecommand{\bcap}{\mathbin{\bk{\cap}}}
\providecommand{\rsm}{\mathbin{\rd{\setminus}}}
\providecommand{\bsm}{\mathbin{\bk{\setminus}}}

% --- comprehension braces of each theory ---------------------------------
\providecommand{\rbr}[1]{\rd{\{}\,#1\,\rd{\}}}
\providecommand{\bbr}[1]{\bk{\{}\,#1\,\bk{\}}}

% --- colour-blind membership (sees through the colour barrier) ------------
\providecommand{\cbin}{\mathrel{\in_{\!\star}}}

% --- FM symbols (shared with rset.stdfmset.tex) --------------------------
\providecommand{\Atoms}{\mathbb{A}}
\providecommand{\Sym}{\mathrm{Sym}}
\providecommand{\pow}{\mathcal{P}}
\providecommand{\supp}{\mathrm{supp}}
\providecommand{\fs}{\mathrm{fs}}
\providecommand{\new}{\mathbin{\reflectbox{\ensuremath{\mathsf{N}}}}}
\providecommand{\atomsof}{\mathrm{atoms}}
\providecommand{\Red}{\rd{\mathsf{R}}}           % the red HF universe
   % red/black colour layer + FM symbols (providecommand-guarded)

\section{A model of $\mathsf{FM}$ in the reflective universe}
\label{sec:fmmodel}

We can now collect the dividend of the construction \emph{Tying the second
recursive knot} above. The reflective theory was assembled from the empty set
alone, yet --- we claim --- it already contains a model of $\mathsf{FM}$, and in
fact two of them, interlocked. The atoms such a model needs are not postulated;
they are \emph{supplied by the dual copy}. Throughout we fix the two tied
theories introduced above, the red copy $\rSetX$ and the black copy $\bSetX$,
with
\[
  \rX \;=\; \bSetX, \qquad \bX \;=\; \rSetX,
\]
so that the admissible atoms of the red theory are exactly the black sets, and
\emph{vice versa}. We write the red copy in {\rd red} and the black copy in
{\bk black} throughout, so that a red comprehension $\rbr{\cdots}$ and a black
comprehension $\bbr{\cdots}$, or the red and black membership relations
$\rin$ and $\bin$, are told apart at a glance.

\subsection{Atoms out of the empty set}
\label{sec:fmmodel:atoms}

Inside the black copy, build the von Neumann numerals from the black empty set,
\[
  \bk{0} \;=\; \bemp,
  \qquad
  \bk{n+1} \;=\; \bk{n} \bcup \bbr{\bk{n}},
\]
each a black, hereditarily finite set obtained, ultimately, from $\bemp$. By
black extensionality (the equations of the formal theory above) these are
pairwise distinct: $\bk{n} \neq \bk{m}$ whenever $n\neq m$. Put
\[
  \Atoms \;:=\; \{\, \bk{n} \;:\; n\in\omega \,\}\ \subseteq\ \bSetX .
\]

Two observations make $\Atoms$ behave \emph{as a set of atoms} for the red
theory.
\begin{itemize}
\item \textbf{Admissibility.} By the second recursive knot $\rX=\bSetX$, every
  black set is an admissible red atom; in particular each $\bk{n}$ is one.
\item \textbf{Opacity.} The red membership relation $\rin$ cannot see into a
  black set (the \emph{Intuitions} above). Hence, \emph{as far as the red theory
  can tell}, an element $\bk{n}$ of $\Atoms$ has no internal structure
  whatsoever: the only facts the red theory may assert about $\bk{n}$ and
  $\bk{m}$ are $\bk{n} = \bk{m}$ or $\bk{n} \neq \bk{m}$, and those are settled,
  out of sight, by black equality.
\end{itemize}
Thus $\Atoms$ presents to the red theory precisely as $\mathsf{FM}$'s set of
atoms does: a countable family of mutually distinct, structureless objects that
may sit inside sets but admit no membership query of their own.

\begin{remark}[Where the risk went]
\label{rem:risk}
This is the sense, promised in the introduction, in which atoms cost no infinite
risk. $\mathsf{ZFA}$ and $\mathsf{FM}$ \emph{postulate} an infinite set of
primitive atoms; the postulate is the risk. Here nothing is postulated beyond
$\bemp$. The atoms are black hereditarily finite sets, and the stream $S[X]$ of
the formal theory hands them out lazily, $\mathrm{take}(\Sigma,n)$ at a time, so
that any red set ever built mentions only finitely many of them. The supply is
\emph{potentially} infinite, never an actual completed infinity: the commitment
is the black empty set together with the same successor step that underwrites
$\omega$, and not one primitive atom more.
\end{remark}

\subsection{The permutation action}
\label{sec:fmmodel:action}

The stream operation $\mathrm{swap}(S[X],\rho)$ of the formal theory applies a
finite permutation $\rho=\{\,\bk{j_1}/\bk{i_1}:\dots:\bk{j_k}/\bk{i_k}\,\}$ to
the ordering of atoms; these finite permutations generate the full group
$\Sym(\Atoms)$. A permutation $\pi\in\Sym(\Atoms)$ acts on red sets by
red-membership recursion exactly as in $\mathsf{FM}$,
\[
  \pi\cdot a = \pi(a)\ (a\in\Atoms),
  \qquad
  \pi\cdot S = \rd{\{}\, \pi\cdot e \;:\; e \rin S \,\rd{\}}.
\]
Because $\pi$ moves only atoms and never the black interior of an atom, and
because the red operations $\remp,\ \rcup,\ \rcap,\ \rsm$ and
$\mathrm{embrace}_i$ are visibly equivariant (a renaming commutes with each), the
action is by automorphisms of the red theory and respects the equations of the
formal theory. We have, verbatim, the $\Sym(\Atoms)$-action that $\mathsf{FM}$
requires.

\subsection{Finite support is automatic}
\label{sec:fmmodel:support}

Write $\Red$ for the class of well-formed red sets, taken modulo red equality,
all of whose atoms lie in $\Atoms$. The finitary rules of the formal theory ---
$\mathrm{embrace}_i$ has finite arity, $\mathrm{Atomicity}$ adjoins one atom, and
$\rcup,\rcap,\rsm$ and $\mathrm{Comprehension}$ over finite sets preserve
finiteness --- yield only \emph{hereditarily finite} red sets. For $S\in\Red$ let
$\atomsof(S)\subseteq\Atoms$ be the finite set of atoms occurring hereditarily in
$S$.

\begin{lemma}[Support]
\label{lem:support}
For every $S\in\Red$ the finite set $\atomsof(S)$ supports $S$: if
$\pi\in\Sym(\Atoms)$ fixes $\atomsof(S)$ pointwise then $\pi\cdot S = S$.
Moreover $\atomsof(S)$ is the \emph{least} support, so $\supp(S)=\atomsof(S)$.
\end{lemma}

\begin{proof}[Proof sketch]
That $\atomsof(S)$ supports $S$ is an induction on the rank of $S$: an atom
$a\in\atomsof(S)$ is fixed by hypothesis, and for a set $S=\rbr{e_1,\dots,e_m}$ a
permutation fixing $\atomsof(S)\supseteq\bigcup_i \atomsof(e_i)$ fixes each $e_i$
by induction, hence fixes $S$. Leastness uses that $\Atoms$ is infinite: if
$a\in\atomsof(S)$ then, choosing a fresh $b\notin\atomsof(S)$, the transposition
$(a\,b)$ fixes $\atomsof(S)\setminus\{a\}$ yet moves $S$, so $a$ lies in every
support. Hence no proper subset of $\atomsof(S)$ supports $S$.
\end{proof}

\begin{corollary}
\label{cor:hfs}
Every $S\in\Red$ is hereditarily finitely supported, with
$\supp(S)=\atomsof(S)$. Freshness is occurrence-failure,
$a\mathrel{\#}S \iff a\notin\atomsof(S)$, and since $\Atoms$ is infinite the
fresh quantifier $\new$ and its some/any law (Section~\ref{sec:fm}) hold in
$\Red$.
\end{corollary}

The point worth stressing is that in $\mathsf{FM}$ finite support is an
\emph{axiom}, imposed by cutting the cumulative hierarchy down to
$\mathcal{U}_{\mathsf{FM}}$. Here it is a \emph{theorem}: the finitary red rules
never build anything that is not finitely supported, so the cut is vacuous and
nothing need be assumed.

\begin{remark}[Two points of precision]
\label{rem:group}
Lemma~\ref{lem:support} rests on two facts worth making explicit. First, the
group in play is the \emph{finitary} symmetric group of $\Atoms$ --- the
permutations generated by the transpositions that $\mathrm{swap}$ realises --- for
which supports are closed under intersection and least supports therefore exist;
on an arbitrary subgroup they need not. Second, the action permutes atoms while
treating each as structureless: it is an automorphism of the \emph{red} structure
$(\Red,\rin)$ alone, and is not asked to respect the black membership $\bin$
internal to an atom. The atoms are accordingly black-distinct but red-opaque, and
the entire support theory lives at the red level.
\end{remark}

\subsection{The representation theorem}
\label{sec:fmmodel:thm}

Let $\mathcal{M}$ be the structure with universe $\Atoms\cup\Red$, with the atoms
$\Atoms$ as urelements, with $\in^{\mathcal M}$ the red membership relation and
$=^{\mathcal M}$ red equality.

\begin{theorem}[Representation]
\label{thm:rep}
$\mathcal{M}$ is a model of hereditarily-finite Fraenkel--Mostowski set theory:
the $\mathsf{FM}$ analogue of $\mathrm{HF}(\Atoms)$. Concretely, interpreting the
atoms of $\mathsf{FM}$ as $\Atoms$, its sets as $\Red$, and its membership as red
$\rin$, the structure $\mathcal{M}$ validates extensionality, the empty set,
pairing, (finite) union, (finite) separation and foundation in their
$\mathsf{ZFA}$ form, carries the $\Sym(\Atoms)$-action of
Section~\ref{sec:fmmodel:action}, and satisfies the finite-support axiom. The
fresh quantifier and the some/any theorem hold throughout.
\end{theorem}

\begin{proof}[Proof sketch]
The atoms are membersless: red $\rin$ cannot enter a black set, so no
$a\in\Atoms$ has red elements, and atoms are $\in$-minimal and disjoint from
$\Red$ in sort, as $\mathsf{ZFA}$ demands. Foundation holds because each red set
is built at a finite rank over the atoms, so $\in^{\mathcal M}$ is well-founded.
The remaining axioms come straight from the rules of the formal theory:
$\mathrm{Foundation}$ gives $\remp$; $\mathrm{embrace}$ with $\mathrm{Union}$
gives pairing $\rbr{e_1}\rcup\rbr{e_2}$ and, iterated, every finite set;
$\mathrm{Union}$ gives binary and hence finite unions; $\mathrm{Subtraction}$ and
$\mathrm{Comprehension}$ over finite domains give finite separation; and
extensionality holds for the model provided provable red equality coincides with
sameness of members --- a property the equations of the formal theory are intended
to force, and which Section~\ref{sec:rigor} records as one of the obligations a
fully rigorous development must discharge. The $\Sym(\Atoms)$-action and the
finite-support axiom are Section~\ref{sec:fmmodel:action} and
Corollary~\ref{cor:hfs}; the some/any law follows because $\Atoms$ is infinite
and every $S\in\Red$ has finite support. Powerset and infinity are absent --- as
expected of an $\mathrm{HF}$ universe --- and are the subject of
Section~\ref{sec:fmmodel:infinitary}.
\end{proof}

\begin{corollary}[Two interlocked models]
\label{cor:two}
By the red/black symmetry of the construction, the same argument run with the
colours exchanged produces a black model $\bk{\mathcal M}$ of
hereditarily-finite $\mathsf{FM}$ whose atoms are the red numerals
$\{\,\rd{n} : n\in\omega\,\}\subseteq\rSetX$. The two models interlock: the atoms
of each are sets of the other. In the game-semantic reading of the
\emph{Intuitions} above, player's atoms are opponent's strategies and conversely
--- a single reflective universe carrying both $\mathsf{FM}$ models at once, each
underwriting the other's names.
\end{corollary}

\begin{remark}[What the $\mathrm{HF}$ restriction costs]
\label{rem:hftrivial}
It should be said plainly what is and is not delivered. The phenomena that give
$\mathsf{FM}$ its character --- finite support as a genuine constraint on
\emph{infinite} objects, equivariance with non-trivial orbits, and the some/any
law quantifying over infinitely many fresh names against a fixed infinite
support --- all require infinite, finitely-supported sets, of which $\mathcal{M}$
has none. On the hereditarily finite fragment these properties hold, but
trivially: every set is finite, so finite support is automatic and the some/any
law degenerates. $\mathcal{M}$ is thus best described as a model of the
\emph{nominal hereditarily finite sets}: faithful to the $\mathsf{ZFA}$ and
$\mathsf{FM}$ axioms, but exercising the distinctively nominal content only once
the infinitary extension of Section~\ref{sec:fmmodel:infinitary} is in force. We
retain the phrase ``model of $\mathsf{FM}$'' only under this qualification.
\end{remark}

\subsection{Finitary by design: the hereditarily finite core}
\label{sec:fmmodel:finitary}

It is worth being explicit about why Theorem~\ref{thm:rep} lands on the
hereditarily finite fragment, because the reason is also the source of the
construction's advantage. The signature of the formal theory offers, as set
constructors, only $\remp$, the finite-arity embraces $\mathrm{embrace}_i$
($|\mathrm{embrace}_i| = i < \omega$), the binary $\rcup,\rcap,\rsm$, and
$\mathrm{Comprehension}$ over finite domains $S_1\times\dots\times S_n$. Each of
these maps finite sets to finite sets, and the only base case is $\remp$;
so a trivial induction shows that \emph{every} derivable red set is hereditarily
finite. The model is $\mathrm{HF}(\Atoms)$, the nominal hereditarily finite sets
over the atoms.

Two of $\mathsf{FM}$'s defining commitments are, in this regime, theorems rather
than axioms.
\begin{itemize}
\item \textbf{Foundation} holds because every red set sits at a finite rank over
  the atoms, so $\in^{\mathcal{M}}$ is well-founded (Theorem~\ref{thm:rep}); we
  never assumed regularity.
\item \textbf{Finite support} holds because a hereditarily finite set mentions
  only finitely many atoms, each of which is its own least support
  (Lemma~\ref{lem:support}); we never assumed the finite-support axiom.
\end{itemize}
Correspondingly there is no infinity axiom and no infinite risk: the stream
$S[X]$ is read only through $\mathrm{take}(\Sigma,n)$ for finite $n$, so at no
point is an actually-infinite set formed. This is the regime in which the
reflective theory strictly improves on axiomatic $\mathsf{FM}$ --- both of the
moves that define $\mathsf{FM}$, foundation and the finite-support cut, are
\emph{derived} from the construction rather than \emph{posited}.

\subsection{Going infinitary, and the return of the cut}
\label{sec:fmmodel:infinitary}

Full $\mathsf{FM}$ also has infinite, finitely-supported sets --- $\Atoms$
itself, or the equivariant powerset $\pow_{\fs}(\Atoms)$ of finitely-supported
sets of atoms --- which the finitary signature cannot build. To reach them one
adds an \emph{infinitary} constructor. The most economical is an $\omega$-ary (or
$\kappa$-ary) embrace that consumes a whole stream rather than a finite prefix,
\[
  \inferrule{\Sigma;\Gamma\vdash S \\ (\sigma,\Sigma') = \mathrm{take}(\Sigma,\omega)}
            {\Sigma';\,\sigma,\Gamma\vdash \mathrm{embrace}_\omega(\sigma)\rcup S}
  \quad\textsc{Infinitary-Embrace},
\]
where $\mathrm{take}(\Sigma,\omega)$ returns the entire stream; equivalently one
adjoins a powerset operator and an infinity axiom and climbs the red cumulative
hierarchy $\mathcal{U}_\alpha(\Atoms)$ through all ordinals. Either way two things
change, and --- this is the point the finitary core conceals --- they change
\emph{together}.

\paragraph{An actual infinity, hence one $\omega$ of risk.}
The lazy stream becomes a completed set; $\Atoms$ itself becomes a red set. By
Remark~\ref{rem:risk} the commitment incurred is exactly black infinity, a single
$\omega$-style axiom. No primitive atoms are reintroduced --- the atoms remain
black hereditarily finite sets built from $\bemp$ --- but the potential infinity
of the finitary core is now actual.

\paragraph{Finite support is no longer free; the cut returns.}
Below $\omega$, finite support was automatic. The infinitary hierarchy is not so
forgiving: it contains sets that have \emph{no} finite support. The canonical
witness is the graph of the stream enumeration,
\[
  E \;=\; \{\,(n,a_n) : n\in\omega\,\},
  \qquad a_n = \text{the $n$-th atom of } S[X],
\]
which is fixed by a permutation $\pi$ only when $\pi$ preserves the enumeration,
i.e. only when $\pi$ is the identity on the infinitely many $a_n$; so no finite
set of atoms supports $E$. This is no accident. The formal theory already tells
us that $S[X]$ ``may be thought of as an ordering on $\Atoms$'', and a total order
on an infinite set of atoms is the textbook example of non-finitely-supported
data. The finitary rules never expose that order --- $\mathrm{take}$ reads only a
finite prefix --- which is precisely why support comes for free below $\omega$.
An infinitary embrace that swallows the whole stream imports the order, and with
it the non-finitely-supported sets.

Consequently, in the infinitary regime one must do exactly what axiomatic
$\mathsf{FM}$ does: cut down to the hereditarily finitely-supported part,
\[
  \mathcal{U}_{\mathsf{FM}} \;=\; \{\, x : x \text{ is hereditarily finitely supported}\,\},
\]
discarding $E$ and its kind. On the cut, the some/any theorem and the
equivariance principle hold verbatim, and the result satisfies the Gabbay--Pitts
axioms. The instructive point is that the cut is now the \emph{only} thing
borrowed from axiomatic $\mathsf{FM}$: the atoms still come from the dual copy,
the $\Sym(\Atoms)$-action still comes from $\mathrm{swap}$, and foundation still
comes from the construction (Section~\ref{sec:nwf}); only the finite-support
axiom must be re-imposed, and only because we elected to admit actual infinities.

This locates precisely where reflection saves risk and where it cannot. Atoms:
always free, from the dual copy. Foundation: always free, from local
well-foundedness (Section~\ref{sec:nwf}). Finite support: free in the finitary
core, an imposed cut once infinities are admitted. The finitary/infinitary
boundary is therefore not a technicality but the exact ledger line between what
the construction \emph{derives} and what it must still \emph{assume} --- and a
deliberate design choice for any theory built this way. We would keep the
finitary signature as the canonical core, where the accounting is cleanest, and
treat \textsc{Infinitary-Embrace} as an optional, clearly-priced extension.
).
%% rset.colors.tex
%% Red/black colour layer for the two intertwined copies, plus the FM symbols
%% used by rset.mainthm.tex and rset.applications.tex.
%%
%% EVERYTHING routes through \rd (red copy = player) and \bk (black copy =
%% opponent): redefine just those two macros to change the whole scheme (e.g.
%% to make the "black" copy dark blue for legibility).  Every definition is
%% \providecommand-guarded, so it is safe to \input this file from several
%% places and harmless if rset.local already defines these names.
%%
%% Requires \usepackage{xcolor} (already loaded in rset.tex).

% --- the two copies -------------------------------------------------------
\providecommand{\rd}[1]{\textcolor{red}{#1}}     % red copy  (player)
\providecommand{\bk}[1]{\textcolor{black}{#1}}   % black copy (opponent)

% --- theory names and atom supplies --------------------------------------
\providecommand{\rSetX}{\rd{\mathrm{Set}[X]}}    % red   Set[X]
\providecommand{\bSetX}{\bk{\mathrm{Set}[X]}}    % black Set[X]
\providecommand{\rX}{\rd{X}}                      % red   atom supply
\providecommand{\bX}{\bk{X}}                      % black atom supply

% --- empty set ------------------------------------------------------------
\providecommand{\remp}{\rd{\emptyset}}
\providecommand{\bemp}{\bk{\emptyset}}

% --- membership (relations) ----------------------------------------------
\providecommand{\rin}{\mathrel{\rd{\in}}}
\providecommand{\bin}{\mathrel{\bk{\in}}}
\providecommand{\rnin}{\mathrel{\rd{\notin}}}
\providecommand{\bnin}{\mathrel{\bk{\notin}}}

% --- binary operations ----------------------------------------------------
\providecommand{\rcup}{\mathbin{\rd{\cup}}}
\providecommand{\bcup}{\mathbin{\bk{\cup}}}
\providecommand{\rcap}{\mathbin{\rd{\cap}}}
\providecommand{\bcap}{\mathbin{\bk{\cap}}}
\providecommand{\rsm}{\mathbin{\rd{\setminus}}}
\providecommand{\bsm}{\mathbin{\bk{\setminus}}}

% --- comprehension braces of each theory ---------------------------------
\providecommand{\rbr}[1]{\rd{\{}\,#1\,\rd{\}}}
\providecommand{\bbr}[1]{\bk{\{}\,#1\,\bk{\}}}

% --- colour-blind membership (sees through the colour barrier) ------------
\providecommand{\cbin}{\mathrel{\in_{\!\star}}}

% --- FM symbols (shared with rset.stdfmset.tex) --------------------------
\providecommand{\Atoms}{\mathbb{A}}
\providecommand{\Sym}{\mathrm{Sym}}
\providecommand{\pow}{\mathcal{P}}
\providecommand{\supp}{\mathrm{supp}}
\providecommand{\fs}{\mathrm{fs}}
\providecommand{\new}{\mathbin{\reflectbox{\ensuremath{\mathsf{N}}}}}
\providecommand{\atomsof}{\mathrm{atoms}}
\providecommand{\Red}{\rd{\mathsf{R}}}           % the red HF universe
   % red/black colour layer + FM symbols (providecommand-guarded)

\section{Applications: nominal techniques founded on reflection}
\label{sec:applications}

The original motivation for $\mathsf{FM}$ in computer science was Pitts and
Gabbay's \emph{nominal} treatment of names and binding
\cite{DBLP:journals/fac/GabbayP02}. That treatment is founded squarely on
$\mathsf{FM}$: names are the atoms, $\alpha$-equivalence is an orbit relation
under the permutation action, and the freshness machinery is exactly support and
the $\new$-quantifier. Since Section~\ref{sec:fmmodel} realises $\mathsf{FM}$
inside the reflective universe, the entire nominal edifice descends with it. The
names that nominal techniques take as primitive become, here, black sets built
from the empty set; the reflective construction \emph{founds} the nominal theory
rather than presupposing its atoms.

\subsection{Nominal sets, internally}
\label{sec:applications:nom}

A \emph{nominal set} is a set carrying a $\Sym(\Atoms)$-action in which every
element is finitely supported; nominal sets and equivariant maps form the
category $\mathbf{Nom}$, equivalent to the Schanuel topos. The model
$\mathcal{M}$ of Theorem~\ref{thm:rep} is, on the nose, a nominal set: the action
is that of Section~\ref{sec:fmmodel:action}, and finite support is
Corollary~\ref{cor:hfs}. Every construction nominal practitioners rely on ---
the equivariance principle (anything definable from equivariant data is
equivariant), least supports, and the some/any law for $\new$ --- is therefore
available inside the reflective universe, with the colour discipline of
Section~\ref{sec:fmmodel} keeping the bookkeeping honest: names live in the
{\bk black} copy, the sets that bind them in the {\rd red} copy.

\subsection{Abstract syntax with binding}
\label{sec:applications:syntax}

The headline application of \cite{DBLP:journals/fac/GabbayP02} is a uniform
account of inductively-defined syntax \emph{modulo} $\alpha$-equivalence. Its key
device is \emph{name abstraction}: for a nominal set $X$,
\[
  [\Atoms]X \;=\; (\Atoms\times X)\big/\!\sim,
  \qquad
  \langle a\rangle x \sim \langle b\rangle y
  \;\;\Longleftrightarrow\;\;
  \new c.\ (a\,c)\cdot x = (b\,c)\cdot y,
\]
where $\langle a\rangle x$ is the class of $(a,x)$ and $(a\,c)$ is a
transposition. Both the action and $\new$ are present in $\mathcal{M}$, so
$[\Atoms]X$ is definable there, and it has the property that makes it a model of
binding: $\supp(\langle a\rangle x) = \supp(x)\setminus\{a\}$, i.e. abstraction
\emph{discharges} the bound name from the support, precisely as a binder removes
a variable from the free-variable set.

Inductive syntax is then obtained as usual. Object-level $\lambda$-terms, for
instance, are the least solution of
\[
  \Lambda \;\cong\; \Atoms \;+\; \Lambda\times\Lambda \;+\; [\Atoms]\Lambda
  \qquad(\text{variable, application, abstraction}),
\]
with $\alpha$-equivalence built into the third summand. Structural recursion and
induction come with the \emph{freshness condition for binders}, justified by the
some/any law: a definition or proof that treats a bound name as fresh for its
parameters is sound because, $\Atoms$ being infinite, ``some fresh name'' and
``every fresh name'' agree (Corollary~\ref{cor:hfs}). This is exactly the
reasoning principle nominal users invoke when they ``choose a fresh variable''.

The risk accounting carries through and sharpens the point. Each individual
term --- each element of $\Lambda$ --- is a hereditarily finite nominal set, so
it already lives in the $\mathrm{HF}$ model $\mathcal{M}$ of
Theorem~\ref{thm:rep}, founded on the empty set alone: a binding tree of any
fixed shape costs no infinite risk. The completed datatype $\Lambda$, an
infinite equivariant set, is one of the finitely-supported infinities of
Section~\ref{sec:fmmodel:infinitary}, and so costs exactly one $\omega$ --- black
infinity --- and still not a single primitive name. Where Pitts and Gabbay begin
with an infinite set of atoms given by fiat, here the names are manufactured, the
$\alpha$-structure is an orbit of the $\mathrm{swap}$ action, and the whole
apparatus is refinanced down to one empty set and one successor.

\subsection{Nominal algebraic theories}
\label{sec:applications:alg}

Beyond syntax, equational reasoning in the presence of names --- nominal
equational logic and the nominal Lawvere theories of Clouston
\cite{DBLP:journals/jcss/Clouston14} --- takes its semantics in $\mathbf{Nom}$.
As $\mathcal{M}$ furnishes $\mathbf{Nom}$-structure internally, any such theory
(name-restriction, scope extrusion, binding signatures and their equational
laws) can be interpreted in the reflective universe. The reflective setting adds
one structural feature these accounts do not have: by
Corollary~\ref{cor:two} there are \emph{two} interlocking name-universes, the red
and the black, each supplying the other's names. Reading red as player and black
as opponent recovers the game-semantic picture of the \emph{Intuitions}: a name
on one side is a strategy on the other, and a nominal theory and its dual coexist
in a single reflective universe.

\subsection{Summary}
\label{sec:applications:summary}

The chain is short and worth stating plainly. Nominal techniques rest on
$\mathsf{FM}$; $\mathsf{FM}$ embeds in the reflective universe
(Section~\ref{sec:fmmodel}); the reflective universe is built from the empty set.
Composing, the Pitts--Gabbay theory of names and binding is founded on two
recursively intertwined copies of ordinary set theory, with its primitive atoms
replaced by constructed black sets and its infinite risk replaced by the finite
risk of $\varnothing$ --- exactly the reconciliation announced in the
introduction, now delivered for the application that made $\mathsf{FM}$ matter to
computer science in the first place.
